% Modernised reading — fonds Alexandre Grothendieck, shelfmark 154
% (pages 1-175, the whole folder).
%
% This is an INTERPRETATION, not a transcription. It re-reads the pages in
% today's notation and vocabulary, states the hypotheses the manuscript leaves
% implicit, and drops the critical apparatus so that the mathematics reads
% straight through. Everything the brackets, underlines, \note{} and \drawing{}
% carried moves into a footnote. It is held to being mathematically correct as
% it stands; where the manuscript is loose or mistaken, the true statement is
% what appears here and the footnote says what the page has. In any dispute of
% substance the transcription governs: whoever cites Grothendieck must cite
% batch-01.fr.tex ... batch-09.fr.tex (pp. 1-20, 21-40, ..., 161-175).
%
% Pass: Opus 5.5 (claude-opus-5-5), 2026-09-27 — first pass, unchecked against the pages by a human.
% The folder was read whole, its nine batches together; this is one document
% for the shelfmark. It was made on Opus 5.5 and has not been measured against
% the reference reading of folder 115 (made on Opus 5). The literature named
% in the footnotes is cited from memory and was not checked against the
% sources. The transcription is a fast first pass with a great many \ill{}:
% much of the connective prose of pp. 3-4, 21-25, 28, 43-51, 85-86 and 89-92
% is not legible, and this reading says so where it matters rather than
% reconstructing it.
%
% Scope. Pages 5, 7 (computer listings) and 10, 98 (separator sheets) are not
% transcribed and are not read here. About half the remaining leaves are
% coloured drawings described in \drawing{} notes; they are summarised in
% footnotes, section by section, not redescribed.
%
% Spine. Given a system Sigma of n pseudo-lines in a real projective plane X,
% build a compact cellular surface whose cells are the relative positions of
% one more pseudo-line D, on which the 2n-gons Pol(D) cut on D by Sigma form a
% local system with dihedral monodromy. Stations (dates in his hand):
%   pp. 1-9     (18.1.84) five variants X_0, X_1, X'_1, X_2, X_3
%   pp. 11-29   (23-25.12.83) the graph Gamma(Sigma) of positions and
%               specialisations, its embedding, weights, polygons; the dual map
%   pp. 30-33   abstract maps: graph + semi-circularisation + transport; flags
%   pp. 33-42   the disc Delta_L, chord systems, transition arcs, gluings
%   pp. 43-54   (31.12.83-2.1.84; p. 47 2.2.84) axioms of the « déploiement »
%               (X, K, L), signs, symmetry bijections, Proposition p. 53
%   pp. 55-63   (3-4.1.84) dihedral monodromy and its characters
%   pp. 64-80   (3.1.84) the eight gluings; the two-week error; true lines
%   pp. 81-100  (5.1.84) surjectivity onto D_2n; generalisation; drawings
%   pp. 101-120 rotation around a face; supersingular band; Euler count
%   pp. 121-149 spindles, repères, sigma, sigma_0, sigma_1, sigma_2; sweeps
%   pp. 150-160 the cellular surface X* = Rel(Sigma); cube; (1.1.84) exceptional
%               facets
%   pp. 161-175 drawings; the axioms again for a surface with boundary (p. 170)
% The leaves are not in date order; the reading follows page order and says
% the dates.
%
% Notation fixed for the folder. Sigma = (D_i)_{i in I}, n = card I, S the
% vertex set, X_1 the union of the D_i. « Supercritique » is used throughout
% for the positions D = D_i (he also says « supersingulière » and
% « exceptionnelle »). Tr(D), Tr*(D) are OUR names for the polygon of
% transitions at D, which pp. 17-24 also call Pol(D), Pol*(D); Pol(D) is kept
% for the 2n-gon of crossings on the double cover D~. \mathbb{D}_{2n} is the
% automorphism group of a combinatorial 2n-gon, of order 4n. At p. 145 the
% adjacent transpositions of a wiring diagram are written s_i (he writes
% sigma_i, which elsewhere names the flag involutions). « Sigma stable »
% (simple arrangement) is distinguished from « D stable ».
%
% Substantive departures, each footnoted where it occurs:
%   - p. 18: « I \ {i_0} » read as the double cover (I \ {i_0})~, as p. 131
%     writes it; the count 2(n-1) requires it;
%   - p. 31: the last line labels sigma_2 sigma_0 as rho_f; it is the edge
%     involution sigma; the face rotation is sigma_1 sigma_0, as the line
%     above and p. 156 have it;
%   - p. 34: the parallelism relation, written « d_i = d_j ou
%     ∂d_i ∩ ∂d_j = ∅ », is read as « d_i and d_j do not cross in Delta »;
%   - p. 45 (i) 4): L = phi^{-1}(X_1) holds only off K_0; the correct
%     statement is phi^{-1}(X_1) = L ∪ K_0, which p. 159 itself gives;
%   - p. 57: D^+_{2n} (the rotations) is Z/2nZ, not Z/nZ;
%   - p. 59-60: the half-turn argument gives surjectivity onto D_2n for n even;
%     for n odd a side reflection is also needed (pp. 63, 85);
%   - p. 62: the boxed chi_X = chi_a chi_c is read chi_X = chi_* chi_c, as the
%     formula before it and the last line of the page give;
%   - p. 85: the rotation subgroup has order 2n; the index-2 subgroup meant is
%     even rotations + side reflections;
%   - p. 119: « chi = 1 - g » is his normalisation (g := 1 - chi), not the
%     genus of a closed surface.
% Flagged, not corrected: the index of a stable exceptional edge, n - (nu_s-1)
% - (nu_t-1) on p. 46, against (n-1) - (nu-1) - (nu'-1) on pp. 55, 93, 110;
% the cancelled table of p. 81 (chi_2 non-trivial on half-turns) against the
% triviality of chi_c claimed on p. 63.
%
% Announced and never established: why there is no local system on X_0
% (p. 6); whether X'_1 maps to X (p. 9); the injectivity of p. 18; the
% programme 1)-2) of pp. 18-19; the specialisation towards a supercritical
% position (pp. 25, 54); the sufficiency of the chord axioms (p. 34); the
% filling of the circuits of exceptional arcs (pp. 39-40); the relation
% between sigma_a^F and the signs (p. 49); the general circuits of A1 (pp. 70,
% 74); the local system without (PS) (p. 92); the pentagonal map (p. 94); the
% value of epsilon (p. 110); the circuits round a hole (p. 112); the
% « rectification b) » of p. 117, which is on no page of the folder; the
% meaning of sigma_0 (pp. 134, 137); the endpoints of an edge a_L of X*
% (p. 150); question 1) of p. 153.
\documentclass[11pt,a4paper]{article}
\input{../preamble/grothendieck}

\folder{154}
\pages{1}{175}
\dating{1983-1984}
\watermark{Édition de démonstration\\interprétation personnelle de l'œuvre}
\foldertitle{[Système de pseudo-droites] : notes manuscrites (1983-1984, s.d.) — lecture modernisée du dossier entier}

\begin{document}

\section*{Résumé}

\begin{resume}
Ces feuillets, écrits entre Noël 1983 et février 1984, cherchent à dessiner la
carte de toutes les façons dont une droite peut se placer parmi d'autres.

Tracez quelques droites sur une feuille, puis posez une règle par-dessus et
faites-la glisser. Tant qu'elle ne passe par aucun des points où les droites se
croisent, rien d'essentiel ne change : elle coupe les mêmes droites, dans le
même ordre. Le « type » de la règle ne change qu'au moment où elle franchit un
point de croisement. Ces types sont ce que le dossier appelle des
\emph{positions relatives}, et il y en a un nombre fini. Deux positions sont
voisines quand on passe de l'une à l'autre en franchissant un seul point. La
question du dossier est : quelle forme a l'ensemble de ces positions, pris
tout entier avec ses voisinages ? Il veut en faire une surface, découpée en
cellules comme une carte de géographie l'est en pays.

Deux choses rendent la question intéressante. D'abord Grothendieck ne travaille
pas avec des droites rigides, mais avec des \emph{pseudo-droites} — des courbes
qu'on peut déformer à volonté, pourvu que deux d'entre elles se coupent toujours
exactement une fois — et dans le plan projectif, où deux droites se coupent
toujours, où il n'y a pas de parallèles et où une droite n'a, en un sens, qu'un
seul côté. Ensuite, la règle emporte avec elle un objet : la suite des points
où elle coupe les autres droites, lue sur ses deux bords, forme un polygone à
$2n$ sommets s'il y a $n$ droites. Quand on promène la règle le long d'un chemin
de positions et qu'on revient au point de départ, ce polygone revient, mais
peut-être tourné ou retourné. C'est le phénomène du pendule de Foucault, ou d'un
vecteur qu'on transporte parallèlement sur une sphère et qui revient d'un autre
angle : la \emph{monodromie}. Le but du dossier est de construire la surface des
positions de telle sorte que ce transport y soit partout bien défini, et de
calculer ce qu'un tour fait subir au polygone.

L'idée qui arrive, et qui organise tout, est de retourner le point de vue :
au lieu de voir chaque position comme un point, la voir comme une petite pièce
de puzzle — le disque qui reste quand on enlève du plan une bande étroite
autour de la règle — et d'obtenir la surface en recollant ces pièces le long de
leurs bords, chaque recollement correspondant au franchissement d'un point de
croisement. Tout le travail est dans le choix de ces recollements : il y en a
huit a priori pour chaque paire de pièces voisines, et seuls certains donnent
une surface sur laquelle le transport du polygone est cohérent. Les droites de
la configuration elles-mêmes, vues comme positions de la règle, posent un
problème à part, et une bonne moitié du dossier cherche comment les faire
entrer dans la surface.

Ce qu'il établit, au bout du compte : la surface des positions existe en
plusieurs variantes ; tourner d'un demi-tour autour de l'une des droites fait
tourner le polygone d'exactement $n - 1$ crans ; et, en combinant ces tours avec
des demi-tours autour des points de croisement, on obtient toutes les
symétries du polygone — le revêtement qui les déploie toutes est d'un seul
tenant. Au passage, il met au net, indépendamment des droites, la manière dont
un graphe muni de petits polygones en chaque sommet définit une surface, par
trois opérations élémentaires sur les « repères » — le formalisme des cartes
qu'il expose au même moment dans l'\emph{Esquisse d'un programme}. Beaucoup
reste ouvert, et le dossier s'arrête sur un dessin de la cellule carrée qu'il
n'a pas fini d'assembler.

On voit ici travailler quelqu'un qui pense en dessinant : la moitié des
feuillets sont des dessins en couleur, et des pages entières ne portent qu'une
droite qui balaie un disque. On le voit aussi se tromper, et le dire. Au début de
janvier, il note : « je m'étais mélangé les pinceaux \emph{deux} fois », et
constate qu'il s'était arrêté pendant deux semaines sur le seul recollement qui
ne marche pas, « Dieu sait pourquoi ! » ; le 18 janvier, devant une question
qu'il ne sait plus trancher : « Je ne suis plus trop dans le coup, et devrait
reconsulter des notes des semaines précédentes. » Et les joies sont notées
aussi franchement : « ce qui est un triomphe ! », « merveilleux, c'est toujours
pair ! », « ouf ! ».

Pour aller plus loin, les noms d'aujourd'hui : arrangements de pseudo-droites
et matroïdes orientés de rang $3$, dont les positions relatives sont les
extensions à un élément, et l'espace de ces extensions ; flips de triangles ;
cartes combinatoires, systèmes de rotation signés et groupe cartographique ;
monodromie et revêtements principaux de groupe diédral ; diagrammes de fils et
suites admissibles.

\keywords{pseudoline arrangement, real projective plane, single-element extension, oriented matroid, extension space, triangle flip, cellular embedding, dual map, crosscap, branched covering, local system, weighted map, combinatorial map, signed rotation system, flag system, cartographic group, barycentric subdivision, Möbius band, chord diagram, folding map, immersion, dihedral group, monodromy, principal bundle, orientation character, projective duality, line arrangement, Euler characteristic, wiring diagram, allowable sequence, reduced word, hypercube}
\end{resume}

\pagerange{1}{175}
\section*{Le fil du dossier, et les conventions}

\subsection*{La donnée}

Tout le dossier part d'une seule donnée. Un \emph{système de pseudo-droites}
est un plan projectif réel topologique $X$ muni d'une famille finie
$\Sigma = (D_i)_{i \in I}$ de pseudo-droites — des courbes fermées simples
qu'on ne peut pas rétracter, dont deux quelconques se coupent en exactement un
point, en s'y croisant. On note $n = \operatorname{card} I$, $X_1 = \bigcup
D_i$ la réunion des pseudo-droites, et $S = S(\Sigma)$ l'ensemble des
\emph{sommets}, points où se coupent au moins deux $D_i$ ; un sommet où passent
$\nu$ pseudo-droites est dit d'ordre, ou de multiplicité, $\nu$. Le système est
dit \emph{simple} quand tous ses sommets sont doubles ; le dossier dit alors
« $\Sigma$ stable » (pages 2, 77, 119), et il faut ne pas confondre cet emploi
avec celui qui suit.\footnote{Le dossier suppose aussi, à l'occasion, que
$\Sigma$ n'est pas \emph{trivial}, c'est-à-dire que ses droites ne passent pas
toutes par un même point (page 11), et il écarte par leur nom quelques
systèmes particuliers : $\mathrm{II}_p$ (page 47), $\mathrm{I}_n$ et
$\mathrm{III}_{p,q}$ (page 96). Ces noms sont ceux de la classification du
dossier 80 ; la page 96 définit elle-même les deux derniers cas par leur
propriété : $\mathrm{III}_{p,q}$ quand toute pseudo-droite passe par l'un des
deux sommets $s$, $t$ considérés, $\mathrm{I}_n$ quand toute passe par $u$ ou
$u'$. Le « $\mathrm{St}_4$ » de la page 13 désigne vraisemblablement
l'arrangement simple de quatre pseudo-droites.}

L'objet de toutes les pages est une pseudo-droite \emph{de plus}, $D$, ajoutée à
$\Sigma$ de telle sorte que $\Sigma \cup \lbrace D \rbrace$ soit encore un
système de pseudo-droites. Sa \emph{position relative} est sa classe à isotopie
près, l'isotopie étant astreinte à ne jamais faire passer $D$ par un nouveau
sommet ni lui en faire quitter un (la page 87 le dit précisément). On note
$P = \mathrm{PR}(\Sigma)$ l'ensemble des positions relatives. La
\emph{spécialité} de $D$ est le nombre $\operatorname{card}(S \cap D)$ des
sommets de $\Sigma$ par lesquels elle passe (page 137). La position est
\emph{stable} si elle ne passe par aucun sommet, \emph{sous-stable} si elle
passe par exactement un (le dossier dit aussi « semi-stable », pages 71, 132,
135, 150), et \emph{critique} si elle en contient au moins deux ; dans les
dernières pages, « critique » veut dire, plus précisément, \emph{minimale pour
la spécialisation} (page 141).\footnote{Aux pages 16 et 153 apparaissent aussi
des sommets « critiques extérieurs » et des positions « principales (donc
critiques) », qui ne sont pas définis ; on les signale où ils figurent.} Les
positions \emph{supercritiques} sont les $D_i$ eux-mêmes, vus comme positions
d'une pseudo-droite relativement aux autres ; le dossier les appelle aussi
« supersingulières » et « exceptionnelles », et on ne garde ici que le premier
mot.\footnote{La page 4 écrit « supercritiques » par-dessus « supersingulières »
biffé ; la graphie du mot reste douteuse en plusieurs endroits de la
transcription. Les pages 57 à 110 disent le plus souvent « supersingulière ».}

On passe d'une position à une autre par deux mouvements inverses. Une
\emph{spécialisation} $D \to D'$ fait glisser $D$ à travers un triangle de
$\Sigma \cup \lbrace D \rbrace$ jusqu'à ce qu'elle passe par un sommet de plus ;
une \emph{générisation} l'en décolle. C'est, dans le vocabulaire d'aujourd'hui,
le « flip » d'un triangle, ou mutation.\footnote{Ce vocabulaire (triangle flip,
mutation) vient de la théorie des arrangements de pseudo-droites et des
matroïdes orientés ; le théorème de Ringel (1956), qui dit que deux
arrangements simples se déduisent l'un de l'autre par de tels mouvements, est
antérieur au dossier, mais rien sur ces pages ne dit que Grothendieck le
connaissait.}

Deux objets attachés à une position reviennent sur toutes les pages. Un
voisinage tubulaire étroit de $D$ dans $X$ est une bande de Möbius ; son bord
est un cercle qui revêt $D$ deux fois, le \emph{revêtement des rives}
$\widetilde{D}$, muni de l'involution sans point fixe qui échange les deux
points au-dessus d'un même point de $D$, l'\emph{antipodisme}. Chaque $D_i$
coupe $D$ une fois, donc $\widetilde{D}$ deux fois : les $2n$ points ainsi
obtenus (comptés avec multiplicité aux sommets) font de $\widetilde{D}$ un
polygone combinatoire d'ordre $2n$, noté $\operatorname{Pol}(D)$. On note
$\omega(D)$ l'ensemble à deux éléments de ses orientations. Le groupe des
automorphismes d'un polygone combinatoire d'ordre $2n$ est le groupe diédral
$\mathbb{D}_{2n}$, d'ordre $4n$ ; son sous-groupe des rotations
$\mathbb{D}^{+}_{2n}$ est cyclique d'ordre $2n$.\footnote{Le dossier écrit
$\mathbb{D}_{2n}$ avec un D ajouré ; nous gardons sa notation, indexée par
l'ordre du polygone et non par celui du groupe. Deux fois (pages 57 et 85), il
attribue au sous-groupe des rotations l'ordre $n$ ; on corrige sur place.} On
note enfin $\widetilde{X}$ le revêtement d'orientation de $X$, une sphère : la
page 28 identifie les points de l'image inverse $\widetilde{\Delta}$ d'une
pseudo-droite $\Delta$ situés au-dessus d'un sommet $s$ aux deux orientations
de $X$ en $s$, c'est-à-dire aux deux rives de $\Delta \smallsetminus \lbrace s
\rbrace$. Les deux revêtements doubles $\widetilde{\Delta}$ (image inverse
dans $\widetilde{X}$, revêtement des rives) sont connexes et donc isomorphes,
et le dossier passe librement de l'un à l'autre.

Deux notations du dossier ont un double sens, et on les sépare. Aux pages 17 à
24, « $\operatorname{Pol}(D)$ » désigne aussi le polygone formé par les
spécialisations et générisations immédiates de $D$, rangées sur
$\widetilde{D}$ ; on l'appelle ici \emph{polygone des transitions} et on le
note $\operatorname{Tr}(D)$, et $\operatorname{Tr}^{*}(D)$ sa version agrandie
des pages 22 à 24 (le « $\operatorname{Pol}^{*}(D)$ » de la page). Et à la page
145, les $\sigma_i$ sont les transpositions de brins voisins d'un diagramme de
fils ; on les note $s_i$, parce que partout ailleurs les $\sigma$ sont les
involutions qui agissent sur les drapeaux d'une carte.

\subsection*{Le but, et deux images duales}

Le but du dossier se dit en une phrase : faire de l'ensemble $P$ des positions
relatives une \emph{surface} compacte, décomposée en cellules, sur laquelle les
polygones $\operatorname{Pol}(D)$ forment un \emph{système local} — de sorte
qu'en suivant un chemin de positions on transporte le polygone de l'une sur
celui de l'autre, et qu'un lacet agisse sur $\operatorname{Pol}(D)$ par un
élément de $\mathbb{D}_{2n}$ (la \emph{monodromie}). Le revêtement principal de
groupe $\mathbb{D}_{2n}$ associé est ce que la page 9 appelle le
« déploiement ultime ».

Deux images de cette surface se succèdent, duales l'une de l'autre, et le
lecteur doit savoir laquelle est en jeu.
\begin{itemize}
\item \textbf{Les positions comme sommets} (pages 11 à 24). Les positions sont
les sommets d'un graphe orienté $\Gamma(\Sigma)$ dont les arêtes sont les
spécialisations immédiates ; ce graphe se plonge dans une surface
$\operatorname{Str}(\Sigma)$.
\item \textbf{Les positions comme faces} (à partir de la page 25 : « c'est la
carte duale qui est la plus intéressante »). Chaque position $D$ est
représentée par un disque $\Delta(D)$, le complémentaire de la bande de Möbius
autour de $D$, et la surface s'obtient en recollant ces disques. C'est
l'image du « déploiement » axiomatisé aux pages 43 et 170, et celle des
pages 121 à 160, où les faces sont les positions stables, les arêtes les
sous-stables et les sommets les critiques.
\end{itemize}
Dans le cas où les $D_i$ sont de vraies droites, cette seconde image est
connue de tous : la page 77 l'identifie à l'arrangement, dans le plan dual
$X^{\vee}$, des droites qui joignent deux à deux les points duaux des $D_i$. En
termes d'aujourd'hui, une position relative est une \emph{extension à un
élément} du matroïde orienté de rang $3$ que définit $\Sigma$, et la surface
cherchée est un modèle cellulaire de son \emph{espace des
extensions}.\footnote{Le rapprochement est le nôtre ; les matroïdes orientés
(Bland, Las Vergnas, Folkman–Lawrence, 1975-1978) et leur théorie des
extensions existaient en 1983, mais rien dans le dossier ne s'y réfère. On sait
aujourd'hui que l'espace des extensions d'un matroïde orienté de rang $3$ est
connexe (Sturmfels et Ziegler, 1993), alors qu'en rang $4$ il peut ne pas
l'être (Mnëv et Richter-Gebert, 1993) — références citées de mémoire.}

\subsection*{Les stations}

Les feuillets sont datés, pour beaucoup, de la main de Grothendieck, et ils ne
sont pas dans l'ordre des dates : la première rédaction (pages 11 à 42) est de
Noël 1983, la mise en axiomes (pages 43 à 54) du 31 décembre et du 2 janvier,
les calculs de monodromie des 3, 4 et 5 janvier 1984, et les cinq pages qui
ouvrent le dossier du 18 janvier ; une page (47) porte le 2 février
1984.\footnote{Dates de sa main : « 23.12.83 » (page 11), « 24.12. » (page 25),
« 25.12 » (page 27), « 31.12.83 » (page 43), « 1.1.84 » (page 159), « 2.1.84 »
(pages 48 et 50), « 3.1.84 » (pages 57 et 64), « 4.1. » (page 60), « 5.1.84 »
(pages 85 et 87), « 18.1.84 » (page 1), « 2.2.84 » (page 47). Celles des pages
1, 11, 47 et 159 ont été vérifiées sur le fac-similé par Michel Hua le 27
septembre 2026 ; la page 11 peut se lire 23 ou 25, la page 159 1.1.83 ou
1.1.84, et la page 47 avait d'abord été lue « 7.1.1984 ». Deux suites
paginées de sa main : 1 à 6 et 6' aux pages 1 à 9, 1 à 14 aux pages 11 à 25.
Les bannières imprimées des listings (août 1982) sont des dates du papier, non
des siennes.} On suit ici l'ordre des pages.
\begin{itemize}
\item \textbf{Cinq variantes de la surface} — pages 1 à 9 : cinq manières de
traiter les positions supercritiques, et leurs défauts.
\item \textbf{La structure de l'ensemble des positions} — pages 11 à 20 : le
graphe $\Gamma(\Sigma)$, ses polygones locaux, les couloirs, la surface
$\operatorname{Str}(\Sigma)$ et les poids.
\item \textbf{Polygones agrandis et carte duale} — pages 21 à 29 :
$\operatorname{Tr}^{*}(D)$, la naissance d'un système local à partir des poids,
puis les disques $\Delta(D)$.
\item \textbf{Les cartes en général} — pages 30 à 33 : un graphe, une
semi-circularisation et un transport des orientations définissent une surface ;
les opérations $\sigma_0$, $\sigma_1$, $\sigma_2$ sur les drapeaux.
\item \textbf{Le disque d'une position et ses arcs} — pages 33 à 42 : systèmes
de cordes, arcs de générisation, de spécialisation, d'inertie ; le recollement,
jusqu'aux positions supercritiques.
\item \textbf{Le déploiement en axiomes} — pages 43 à 54 : le triplet
$(\mathcal{X}, K, L)$, l'application $\varphi : \mathcal{X} \to X$, les signes
$\varepsilon_F(a)$, les bijections de symétrie, la Proposition de la page 53.
\item \textbf{La monodromie diédrale} — pages 55 à 63 : le transport le long
d'un demi-tour, les images de $\pi_1$ dans $\mathbb{D}_{2n}$, les caractères.
\item \textbf{Huit recollements} — pages 64 à 80 : le choix du recollement, une
erreur de deux semaines, la surface des positions stables, le cas des vraies
droites.
\item \textbf{Surjectivité et généralisation} — pages 81 à 100 : le revêtement
principal est connexe ; les positions d'un cercle par rapport à un graphe sur
une surface quelconque.
\item \textbf{Rotation autour d'une face} — pages 101 à 120 : l'opération
$\rho(D, \omega, a)$, la bande supercritique, la caractéristique d'Euler.
\item \textbf{Fuseaux et repères} — pages 121 à 149 : les quatre involutions
$\sigma_1$, $\sigma$, $\sigma_2$, $\sigma_0$ ; les balayages d'un disque de sept
pseudo-droites.
\item \textbf{La surface $X^{*}$ des positions} — pages 150 à 160 : sa structure
locale, les types de fuseaux, le cube des générisations ; les facettes
exceptionnelles.
\item \textbf{Dessins, et les axiomes repris} — pages 161 à 175.
\end{itemize}

Deux reprises ne doivent pas être confondues avec des répétitions : les axiomes
du déploiement sont écrits deux fois, pour une surface fermée (page 43) et pour
une surface à bord (page 170) ; et le choix du recollement de deux disques est
fait une première fois aux pages 37-38, puis refait entièrement aux pages 64 à
70, où Grothendieck constate qu'il s'était trompé. La moitié des feuillets, à
partir de la page 80 surtout, sont des dessins en couleur ; ils ne sont pas
redessinés ici, et chaque section dit en note ce qu'ils montrent.

\subsection*{Ce qui est annoncé et ne l'est pas}

Beaucoup de choses sont posées comme questions et laissées là, et ce dossier est
un état de travail : les sections le disent à mesure. Les principales :
pourquoi les polygones ne forment pas de système local sur $\mathcal{X}_0$
(page 6) ; l'injectivité de l'application de la page 18 ; la suffisance des
axiomes de cordes (page 34) ; le bouchage des circuits d'arcs exceptionnels
(pages 39-40) ; la relation entre bijections de symétrie et signes (page 49) ;
les circuits généraux de la variante $A_1$ (pages 70 et 74) ; la carte
pentagonale (page 94) ; la valeur du signe $\varepsilon$ (page 110) ; le sens
de l'involution $\sigma_0$ (pages 134 et 137) ; les extrémités d'une arête de
$X^{*}$ (page 150) ; la question 1) de la page 153. La « rectification b) » à
laquelle renvoie la page 117 ne figure sur aucune page du dossier.

\pagerange{1}{9}
\section*{Cinq variantes de la surface de déploiement (pages 1 à 9, 18 janvier 1984)}

Ces pages sont les dernières écrites de la première moitié du dossier, et elles
en font le bilan : il existe au moins cinq façons raisonnables de fabriquer la
surface des positions relatives, qui ne diffèrent que par la manière de traiter
les positions \emph{supercritiques}, et chacune a son défaut.\footnote{Titre de
sa main : « Conventions particulières et constructions relatives aux pos.\ rel.\
supersingulières » ; en tête de la colonne de droite, encadré : « Les 5
variantes de la surface de déploiement d'un syst.\ de pseudo-droites ». Les
pages 3, 4 et 6, d'une écriture rapide, sont en grande partie illisibles ; ce
qui suit en retient ce que les phrases lisibles établissent, et les notes disent
où la lecture manque.} Dans toutes, les faces correspondent aux positions
relatives et les arêtes aux spécialisations — c'est la seconde des deux images
de la section précédente —, et l'on note $F_D$ la face d'une position $D$.

\subsection*{$\mathcal{X}_0$ : les positions supercritiques exclues}

On n'admet pas les positions supercritiques parmi les faces. Elles deviennent
alors des \emph{sommets} de la surface $\mathcal{X}_0$, les « sommets
supercritiques », et ce sont les seuls qui ne soient pas d'ordre $4$ : ils sont
d'ordre $4n$.\footnote{C'est le nombre que porte la page. Dans le cas des vraies
droites (page 77), le point dual $\delta_i$ d'une droite $D_i$ d'un système
simple est un sommet où passent les $n - 1$ droites $\delta_i\delta_j$, donc
d'ordre $2(n-1)$ dans l'arrangement dual ; nous ne savons pas réconcilier ce
compte avec le « $4n$ » de la page, qui compte peut-être autre chose.}

\subsection*{$\mathcal{X}_1$ : une face par position supercritique}

On admet au contraire chaque position supercritique $D_i$ comme une face $F_i$,
et les spécialisations d'une position sous-stable vers elle comme des arêtes,
correspondant à des \emph{arcs de spécialisation} exceptionnels qu'on choisit
aussi courts que possible — de longueur combinatoire $1$ — et placés à égale
distance des arcs de générisation les plus proches. Autour de $F_i$ se
raccordent ainsi, le long de $2n_i$ arcs, les faces des positions sous-stables
voisines, où $n_i$ est le nombre de sommets de $\Sigma$ sur $D_i$ ($2(n-1)$
arcs quand $\Sigma$ est simple) ; les faces des positions stables voisines
viennent s'intercaler, chacune entre deux faces sous-stables adjacentes.

Tous les sommets de $\mathcal{X}_1$ sont alors d'ordre $4$, mais le recollement
laisse des \emph{bords libres}. Pour $\Sigma$ simple, leur nombre dépend de la
parité de $n$ : si $n$ est pair, chaque face sous-stable en donne deux et chaque
face stable qui touche $F_i$ un de plus, soit
\[
2 \cdot 2(n-1) + 2(n-1) = 6(n-1),
\]
« ce qui fait un peu beaucoup » ; si $n$ est impair, chaque face sous-stable en
donne un, au milieu de l'arête commune, et les faces stables aucun, soit
$2(n-1)$, exactement ce qu'il faut pour les recollements.\footnote{Ces comptes
sont ceux de la page 2, que nous n'avons pas refaits ; la fin de l'alinéa b),
qui dit à quoi servent ces $2(n-1)$ bords, est en partie illisible.}

\subsection*{$\mathcal{X}'_1$ : un petit disque à la place du sommet}

Une manière de s'en tirer dans tous les cas part de $\mathcal{X}_0$ : on
remplace chaque sommet supercritique par un petit disque. Chaque face incidente
contribue alors un bord libre, soit $2(n-1)$ bords qui entourent le disque, et
l'on peut les raccorder de façon à reconstituer « l'affinisation
supercritique » de $\Sigma$ sur $F_i$ — l'image affine qu'on obtient en
envoyant $D_i$ à l'infini. Sur la surface $\mathcal{X}'_1$, les sommets situés
sur les faces supercritiques sont d'ordre $3$ et tous les autres d'ordre $4$. Il
faut alors des arcs de spécialisation de longueur combinatoire quelconque, et
passer à la subdivision barycentrique.\footnote{Figures des pages 1 à 3 : un
faisceau de droites au crayon passant par un sommet $s$, légendé
$\mathcal{X}_0$, « $(n = 5)$ » ; la face supercritique $F_i$ en jaune, avec les
faces sous-stables et stables qui s'y raccordent, légendée $\mathcal{X}_1$ ; un
petit disque traversé de pseudo-droites d'où rayonnent des traits rouges et
jaunes, légendé $\mathcal{X}'_1$ pour $n = 5$.}

\subsection*{Le défaut commun}

Sur ces trois surfaces, le système local des polygones $\operatorname{Pol}(D)$,
défini par transport parallèle, n'existe que sur le complémentaire des sommets
ou des faces supercritiques, qui sont des lieux de
\emph{ramification}.\footnote{La lecture de « ramification » est douteuse ; le
reste de l'alinéa, où il énumère un second défaut qui concerne l'application
dans $X$, est illisible sur plusieurs lignes.} Le revêtement principal associé
est donc, pour elles, un revêtement ramifié.

\subsection*{$\mathcal{X}_2$ : éclater le sommet}

La quatrième surface part encore de $\mathcal{X}_0$, mais « fait éclater » les
sommets supercritiques en ajoutant des arcs de transition supplémentaires, de
poids $2$. On l'obtient aussi à partir de $\mathcal{X}'_1$ en découpant
l'intérieur des faces supercritiques et en identifiant les points antipodiques
de leurs bords — c'est-à-dire en remplaçant chaque face supercritique par un
cross-cap, une bande de Möbius. Tous les sommets sont d'ordre $4$, et cette
fois \emph{les $\operatorname{Pol}(D)$ forment un système local}, alors que sur
$\mathcal{X}_0$ il est rompu autour de chaque sommet supercritique. En
revanche $\mathcal{X}_2$ ne s'envoie pas dans $X$ (« sauf erreur »).

En face de cette phrase, entourée d'un grand cercle rouge, une remarque qui
dit l'état du travail :
\begin{quote}
Mais pour quelles raisons, au juste, ne serait-il pas vrai qu'on a un syst.\
local sur $\mathcal{X}_0$ (ce qui est pourtant faux !) Il y a un pt qui en ce
moment m'échappe. Je ne suis plus trop dans le coup, et devrait reconsulter des
notes des semaines précédentes.
\end{quote}
La réponse n'est pas écrite. Celle que suggèrent les calculs des pages 57 à
59, écrits quinze jours plus tôt, est la suivante, et elle est la
nôtre.\footnote{Nous la donnons parce que le dossier en contient tous les
éléments, non parce qu'il la formule. Elle suppose qu'un petit lacet autour du
sommet supercritique de $\mathcal{X}_0$ correspond au « tour complet » autour
de $D_i$ de la page 59, ce que la page 1 ne dit pas.} Un système local sur le
complémentaire d'un point s'étend à travers ce point si et seulement si sa
monodromie autour du point est triviale. Or un tour complet autour d'une
position supercritique agit sur $\operatorname{Pol}(D)$ par le carré de la
translation $t \mapsto t - (n-1)$ de la page 57, c'est-à-dire par
$t \mapsto t - 2(n-1)$ sur $\mathbb{Z}/2n\mathbb{Z}$, qui n'est l'identité que
pour $n = 1$. Le cross-cap de $\mathcal{X}_2$ résout exactement cette
difficulté : son âme réalise un demi-tour, dont le carré est le lacet du bord.

\subsection*{$\mathcal{X}_3$ : une bande de Möbius à la place de $D_i$}

La dernière variante garde les positions supercritiques parmi les faces, mais
en distinguant, pour une $D_i$ portant $n_i$ sommets, $2n_i$ positions
supercritiques, suivant la facette de $\widetilde{D}_i$ où tombe le point
d'intersection initial avec la position voisine. Tous les sommets sont encore
d'ordre $4$ ; la pseudo-droite $D_i$ est remplacée par une bande de Möbius. On
l'obtient à partir de $\mathcal{X}'_1$ en découpant les faces supercritiques
et en recollant, au lieu d'identifier les points antipodiques, une bande dont
les « faces de bord » sont des positions voisines de $D_i$. Cette fois on a une
application $\mathcal{X}_3 \to X$ telle que l'image inverse de $X_1$ soit ce
qu'il faut, $D_i$ compris. Mais on perd la représentation affine de $\Sigma$
qu'offrait $\mathcal{X}'_1$ — la surface représente plutôt un voisinage
tubulaire de $D_i$ — et $\mathcal{X}_3$ ne se prête pas à une interprétation
par la combinatoire interne du système local des $\operatorname{Pol}(D)$.

\subsection*{Les déploiements ultimes}

À chacune des cinq surfaces est associé le revêtement principal de groupe
$\mathbb{D}_{2n}$ défini par les $\operatorname{Pol}(D)$ — ramifié pour les
trois premières. Ce sont ces revêtements qui sont les « déploiements
ultimes ».\footnote{La page les note $\widetilde{\mathcal{X}}_i$. Nous ne leur
donnons pas de symbole, parce que le tilde sert, à partir de la page 57, à une
autre surface : celle où l'on s'interdit de franchir les positions
supercritiques.} Les plus utiles lui semblent $\mathcal{X}'_1$,
$\mathcal{X}_2$, $\mathcal{X}_3$ et leurs revêtements ; peut-être peut-on se
contenter de $\mathcal{X}'_1$, pourvu qu'elle s'envoie bien dans $X$ — ce dont
il n'est « pas tout à fait convaincu », faute de voir ce qui se passe au
voisinage du bord d'une face supercritique.\footnote{Page 9, numérotée 6' :
complément de la page 8. Le reste de la phrase est illisible.}

\pagerange{11}{20}
\section*{La structure de l'ensemble des positions relatives (pages 11 à 20, 23 décembre 1983)}

C'est le début chronologique du dossier, une suite numérotée de 1 à 14 qui court
jusqu'à la page 25.\footnote{Titre de sa main : « Structure sur l'ens.\ des
positions relatives $P = \mathrm{PR}(\Sigma)$ d'un syst.\ $\Sigma = (X, \Sigma)$
de pseudodroites », le $\Sigma$ du système étant souligné deux fois sur la
page.} Les positions y sont les \emph{sommets} d'un graphe.

\subsection*{Les flèches de spécialisation}

Soit $D$ une position relative. Un \emph{triangle de spécialisation} pour $D$ est
la donnée d'un arc de $\widetilde{D}$ formé de $\nu \geqslant 1$ arêtes
consécutives de $\operatorname{Pol}(D)$, dont les $\nu + 1$ extrémités ne sont
pas des sommets de $\Sigma$, découpées par $\nu + 1$ pseudo-droites de $\Sigma$
qui concourent en un même sommet $s$, d'ordre exactement $\nu + 1$, le triangle
de sommets $s$ et des deux extrémités de l'arc ne contenant aucun autre sommet
de $\Sigma$.\footnote{La définition est dans un passage encadré et barré de
quatre traits obliques, qui se lit malgré tout ; rien d'autre sur la page ne
définit la notion, et la suite en a besoin. On exclut le cas où $\Sigma$ est
trivial, « qu'il vaut mieux traiter à part ».} En faisant glisser $D$ à travers
ce triangle jusqu'à $s$, on obtient une position $D'$, sa \emph{spécialisation}
au-dessus de $T$. Soit $\widetilde{P}$ l'ensemble des couples $(D, T)$.

\textbf{L'application $\widetilde{P} \to P \times P$, $(D, T) \mapsto (D, D')$,
est injective.} Car $D' \cap S = (D \cap S) \cup \lbrace s \rbrace$, ce qui
détermine $s$ ; aucune pseudo-droite de $\Sigma \cup \lbrace D \rbrace$ ne
bordant un bigone (sinon $\Sigma$ serait trivial de centre $s$), les arêtes qui
joignent $s$ aux points de $D$ à l'intérieur de $T$ sont uniques, et $D \cap T$
ne contient pas d'autre sommet : $T$ est déterminé. On obtient donc un
\emph{graphe orienté} $\Gamma(\Sigma)$, d'ensemble de sommets $P$ et
d'ensemble d'arêtes $\widetilde{P}$.

Pour $D \in P$, soit $\operatorname{Sp}_0(D)$ l'ensemble de ses spécialisations
immédiates — en bijection avec ses triangles de spécialisation — et
$\operatorname{Gn}_0(D)$ celui de ses générisations immédiates, en bijection
avec les couples $(s, \tilde{\xi})$ où $s \in D \cap S$ et $\tilde{\xi}$ est
l'un des deux points de $\widetilde{D}$ au-dessus de $s$ : décoller $D$ de $s$
d'un côté ou de l'autre. Chaque spécialisation définit un arc de
$\widetilde{D}$, ces arcs sont deux à deux disjoints, et chaque générisation un
point de $\widetilde{D}$, distinct des autres et hors de ces arcs. D'où une
\emph{structure polygonale} — un ordre cyclique défini au sens près — sur
\[
\operatorname{Sp}_0(D) \amalg \operatorname{Gn}_0(D),
\]
l'ensemble des arêtes de $\Gamma(\Sigma)$ issues du sommet $D$ : c'est le
\emph{polygone des transitions} $\operatorname{Tr}(D)$.\footnote{Page 12. La
page ne donne pas de nom à ce polygone ; aux pages 17 à 24 elle l'appelle
$\operatorname{Pol}(D)$, d'où le nom que nous lui donnons (voir les
conventions).}

\subsection*{Au moins trois transitions}

Ce polygone a-t-il au moins trois éléments ? Au moins quatre si $D$ contient
deux sommets de $\Sigma$, par les seules générisations. Pour $D$ stable, il y a
au moins trois triangles de $\Sigma \cup \lbrace D \rbrace$ adjacents à $D$ ;
l'un est contenu dans un unique triangle de spécialisation $T$, et chacun des
deux grands triangles que $T$ découpe de part et d'autre en contient un autre,
d'où trois triangles de spécialisation distincts. Pour $D$ sous-stable il
faudrait au moins un triangle de spécialisation, et l'argument tenté échoue :
la page 13 donne elle-même un contre-exemple, dans l'arrangement simple de
quatre pseudo-droites, avec une position sous-stable qui évite les deux
quadrilatères incidents à son sommet.\footnote{« Il y a effectivement un
contre-exemple avec $\Sigma = \mathrm{St}_4$ » ; les derniers mots de la phrase
sont illisibles. L'énoncé « au moins un triangle de spécialisation pour $D$
sous-stable » est donc faux tel quel ; ce qui suit le rend vrai en comptant les
couloirs.}

La raison en est aussitôt trouvée : « j'ai oublié les cas les plus importants :
les p.r.\ \emph{supercritiques} ». Une position sous-stable $D$, de sommet $a$,
peut se spécialiser vers une pseudo-droite $\Delta$ de $\Sigma$ elle-même, à
condition que $a$ soit un sommet de $\Delta$ et que l'un des deux secteurs
biangulaires de $X$ délimités par $D$ et $\Delta$ soit un
\emph{couloir}\footnote{Le mot n'est défini nulle part. Les figures des pages 13
et 117 montrent une bande entre $D$ et $\Delta$ que les autres pseudo-droites
traversent d'un bord à l'autre, sans sommet à l'intérieur ; c'est le sens
qu'on lui prête ici.} ; la spécialisation se
fait dans ce couloir, le \emph{couloir de spécialisation}, qui compte comme
l'équivalent d'un triangle de spécialisation. Relevé dans $\widetilde{D}$, un
couloir est un arc ouvert d'extrémités les deux points $a'$, $a''$ au-dessus de
$a$, et les spécialisations de $D$ sont alors décrites par des arcs ouverts
deux à deux disjoints de $\widetilde{D} \smallsetminus \lbrace a', a''
\rbrace$. Ainsi complété, le polygone $\operatorname{Tr}(D)$ a toujours au moins
trois éléments, et au moins quatre si $D$ n'est pas stable (page 15, affirmé
sans autre démonstration).

Vues depuis la position supercritique $\Delta$, ses générisations immédiates
sont de deux sortes : deux positions stables par arête de $\Delta$, qui la
longent de part et d'autre, et deux positions sous-stables par sommet,
correspondant aux deux orientations de $X$ en ce sommet. Elles sont donc en
bijection avec les sommets et les arêtes du polygone topologique
$\widetilde{\Delta}$, ce qui fait un polygone d'ordre $4\nu$ si $\nu$ est le
nombre de sommets de $\Delta$, ou d'ordre $2\nu$ si l'on se borne aux
générisations immédiates. L'identification avec le polygone d'une position
ordinaire voisine dépend du choix d'une orientation de $\Delta$.\footnote{Page
14, dont une grande partie se lit mal ; la note de sa main au bas de la page
dit : « ou, d'ordre $2\nu$ seulement en se bornant aux gén.\ \emph{immédiates} ».
La page constate aussi que les couloirs ne forment pas, avec les autres
transitions, une structure polygonale naturelle, ce qui motive ce changement de
point de vue.}

\subsection*{Orientations, plongement cellulaire, poids}

Chaque sommet $D$ de $\Gamma(\Sigma)$, positions supercritiques comprises, porte
donc un polygone $\operatorname{Tr}(D)$ ; soit $\omega(D)$ l'ensemble de ses
deux orientations, qu'on identifie à celui des orientations de $D$. Une flèche
de spécialisation $D \to D'$ définit une bijection $\omega(D) \to \omega(D')$,
et une injection $\operatorname{Gn}_0(D) \hookrightarrow
\operatorname{Gn}_0(D')$ dont le complémentaire a deux éléments.\footnote{Ce
dernier point est la note \emph{NB} qui court de la page 15 à la page 16 ; elle
identifie $\operatorname{Gn}_0(D)$ à l'image inverse de $D \cap S$ dans
$\widetilde{D}$.}

Ces données — un ordre cyclique au signe près en chaque sommet, et une bijection
entre les orientations le long de chaque arête — définissent un
\emph{plongement cellulaire canonique} de $\Gamma(\Sigma)$ dans une surface
compacte, que la page 16 appelle la \emph{surface structurante}
$\operatorname{Str}(\Sigma)$ ; la construction générale est celle des pages 30 à
33. Autour d'un sommet stable, toutes les arêtes partent ; autour d'un sommet
sous-stable, exactement deux arrivent ; autour d'un sommet critique, au moins
quatre arrivent ; et autour des sommets que la page appelle « critiques
extérieurs », toutes arrivent.\footnote{Ce sont vraisemblablement les positions
supercritiques, dont toutes les transitions sont des générisations (page 17) ;
la page ne le dit pas.}

Chaque arête reçoit un \emph{poids}, entier $\geqslant 1$ : le nombre d'arêtes
de $\operatorname{Pol}(D)$ qu'intercepte l'arc de spécialisation. Chaque secteur
angulaire en $D$ — un arc de $\widetilde{D}$ compris entre deux transitions
consécutives — reçoit comme poids le nombre d'arêtes qu'il contient. On obtient
une \emph{carte cellulaire à graphe orienté, à arêtes et à coins pondérés}, et
en chaque sommet non supercritique
\[
\sum \text{poids des arêtes en } D \;+\; \sum \text{poids des coins en } D
\;=\; 2n,
\]
puisque ces arcs partagent les $2n$ arêtes de $\operatorname{Pol}(D)$.

\subsection*{Le polygone déployé, et l'ensemble $\widetilde{I}$}

En éclatant chaque sommet de $\operatorname{Tr}(D)$ en autant de sommets que
l'indique son poids, et en plaçant sur chaque arête autant de points que
l'indique le poids du coin correspondant, on reconstitue le polygone de
$\widetilde{D}$ : on obtient un polygone $\widehat{\operatorname{Pol}}(D)$
dont les sommets correspondent aux arêtes de $\operatorname{Pol}(D)$ et les
arêtes aux points où les $D_i$ coupent $\widetilde{D}$ — le polygone dual de
$\operatorname{Pol}(D)$, visualisé comme un contour parallèle à $D$. Soit
$\widetilde{I}$ l'ensemble des pseudo-droites orientées $\vec{D}_i$, revêtement
double de $I$. On a une bijection canonique
\[
\widetilde{I} \;\simeq\; \lbrace \text{arêtes de }
\widehat{\operatorname{Pol}}(D) \rbrace ,
\]
compatible aux antipodismes.

Pour $D = \Delta = D_{i_0}$ supercritique, il n'y a que des arêtes de
générisation ; on affecte à celle qui correspond au sommet $s$ de $\Delta$ le
nombre de pseudo-droites autres que $\Delta$ qui passent par $s$, et l'on
obtient de même un polygone $\widehat{\operatorname{Pol}}(\Delta)$ d'ordre
$2(n-1)$ — celui que découpent les $D_i$, $i \neq i_0$, sur une parallèle à
$\Delta$ — et une bijection
\[
\widetilde{I \smallsetminus \lbrace i_0 \rbrace} \;\simeq\; \lbrace
\text{sommets de } \widehat{\operatorname{Pol}}(\Delta) \rbrace ,
\]
compatible aux antipodismes, qui dépend du choix d'une orientation de
$\Delta$ : en la changeant, on la compose avec l'antipodisme.\footnote{La page
18 écrit $I \smallsetminus \lbrace i_0 \rbrace$ sans tilde ; un polygone d'ordre
$2(n-1)$ a $2(n-1)$ sommets, et il faut donc le revêtement double, comme
l'écrit d'ailleurs la page 131 : « $\operatorname{Circ}(\widetilde{I \setminus
\lbrace i_0 \rbrace})$ ».} Le poids ordinaire $\nu_0$ d'une arête issue de
$\Delta$ — le nombre d'arêtes découpées par $\Sigma$ sur la générisation le long
du couloir — et le poids $\nu$ qu'on vient de définir vérifient
\[
\nu + \nu_0 = n .
\]

On a ainsi une application canonique
\[
P \smallsetminus \Sigma \;\longrightarrow\; \lbrace \text{structures
polygonales sur } \widetilde{I} \text{ compatibles à l'antipodisme} \rbrace ,
\]
des positions non supercritiques vers les ordres cycliques de $\widetilde{I}$,
dont « il ne devrait pas être difficile de montrer » qu'elle est
injective.\footnote{Rien, dans le dossier, ne le montre. En termes d'aujourd'hui,
cette application associe à une position l'ordre dans lequel elle rencontre les
pseudo-droites orientées ; nous ne savons pas si cet ordre la détermine.} Il
resterait, pour bien faire, à voir (1) comment se raccordent les polygones
$\widehat{\operatorname{Pol}}$ le long d'une flèche de spécialisation et (2)
comment lire toute la structure de graphe orienté — plongement, poids des arêtes
et des coins — sur ces structures polygonales.

Pour (1), la page 19 constate que $\operatorname{Pol}(D)$ et
$\operatorname{Pol}(D')$ « sont les mêmes », sauf au voisinage des deux segments
antipodiques qui correspondent à l'extrémité de l'arête de spécialisation, et
qu'il existe un \emph{isomorphisme canonique} $\widehat{\operatorname{Pol}}(D)
\simeq \widehat{\operatorname{Pol}}(D')$ qui envoie l'un des sommets issus de
cette arête sur l'\emph{antipodique} du sommet correspondant de
$\widehat{\operatorname{Pol}}(D')$, compatible aux poids.\footnote{Le milieu de
la colonne de droite de la page 19, où cette compatibilité est détaillée, est
biffé sur plusieurs lignes et ne se lit qu'en partie. La page 20 est un dessin
en couleur de la même situation, sans texte.}

\pagerange{21}{29}
\section*{Polygones agrandis, système local, carte duale (pages 21 à 29, 23 au 25 décembre 1983)}

\subsection*{Le polygone $\operatorname{Tr}^{*}(D)$}

Les pages 21 à 24 sont d'une écriture très rapide et se lisent mal ; leur
mouvement est cependant clair.\footnote{Une grande partie des pages 21 à 25 est
faite de mots illisibles ; les énoncés retenus ici sont ceux qui se lisent en
entier. Figures : des faisceaux de pseudo-droites sur une droite verte, les arcs
découpés marqués de traits noirs (pages 21, 22, 23) ; la page 26 est une page
de dessins sans texte.} On agrandit $\operatorname{Tr}(D)$ en y ajoutant les
antipodiques des triangles de spécialisation — des
« pseudo-spécialisations » — et des bords libres : c'est le polygone
$\operatorname{Tr}^{*}(D)$, stable par antipodisme. Le poids d'une arête de
$\operatorname{Tr}^{*}(D)$ peut être nul, entre deux générisations. Une
spécialisation $D \to D'$ induit une application $\operatorname{Tr}^{*}(D) \to
\operatorname{Tr}^{*}(D')$ qui n'est pas en général un isomorphisme, parce
qu'une spécialisation peut créer des pseudo-spécialisations lorsque l'autre face
adjacente au triangle $T$ est un couloir. Si $b$ est le sommet qui disparaît,
entre les sommets $x$ et $a$, les poids sont additifs :
\[
\text{poids de } \widehat{ax} \text{ dans } \operatorname{Tr}^{*}(D') \;=\;
\text{poids de } \widehat{xb} + \text{poids de } b + \text{poids de }
\widehat{ba} \text{ dans } \operatorname{Tr}^{*}(D).
\]
C'est la relation de la page 23.\footnote{La page porte
$\operatorname{Pol}^{*}(D)$ dans le premier membre ; la phrase qui précède
demande $D'$, et c'est la lecture de la transcription. L'exposant de
$\operatorname{Pol}^{*}$ est tantôt une étoile, tantôt un signe proche de
$\alpha$. Les types de sommets de $\operatorname{Tr}^{*}(D)$ que la page 23
énumère ensuite, $\alpha$), b), c), sont en majeure partie illisibles.}

La structure de graphe orienté, cellulaire et pondéré, de $\Gamma^{*}(\Sigma)$
ne semble pas distinguer, parmi les arêtes issues d'un sommet, celles qui
correspondent à des pseudo-spécialisations. Mais elle les distingue dès qu'on
connaît les applications $\operatorname{Tr}^{*}(D) \to \operatorname{Tr}^{*}(D')$,
et celles-ci se déduisent de la structure — « \emph{Mais bien sûr} » — en
passant aux polygones déployés : il suffit de savoir comment se correspondent
les orientations, ce qui est connu, et quelle est l'image d'un sommet
particulier $(D, a)$, qui est l'antipodique, dans
$\widehat{\operatorname{Pol}}(D')$, du sommet $(D', a)$.

\subsection*{Un principe général}

La page 24 en tire un principe qui vaut pour toute carte, et c'est le premier
énoncé propre du dossier. \emph{Soit une carte dont les arêtes et les coins
sont pondérés par des entiers $\geqslant 0$, de sorte qu'en chaque sommet la
somme des poids des arêtes et des coins incidents soit un même entier $N$.
Éclatons chaque sommet en un polygone d'ordre $N$, dont les côtés sont répartis
entre les arêtes et les coins selon leurs poids. Alors, à travers chaque arête,
les polygones de ses deux extrémités se correspondent, et ces transports forment
un système local de polygones d'ordre $N$ sur la surface sous-jacente à la
carte.} Pour les positions relatives, $N = 2n$, et le système local est celui
des $\operatorname{Pol}(D)$.\footnote{La page dit : « chaque fois qu'il y a une
2-cellule (avec bords isolés i.e.\ arêtes libres) dont les (vrais) sommets sont
d'ordre $2$, et dont les arêtes et coins sont pondérés (\ldots), la pondération
totale de chaque sommet (\ldots) étant fixe, [on] trouve des (\ldots) de
transport entre les polygones locaux déployés, qui forment donc un syst.\ local
sur la surface sous-jacente à la carte ! ». La formulation que nous en donnons
est la nôtre ; ce qu'elle suppose de plus — que la correspondance à travers une
arête est bien définie par les poids et les orientations — est ce que les pages
21 à 24 établissent dans le cas des positions.}

\subsection*{La carte duale : les disques $\Delta(D)$ (24 décembre)}

Le 24 décembre, changement de point de vue : « c'est la carte \emph{duale}
qui est la plus intéressante ». Le polygone de $D$ se voit comme le bord d'un
petit voisinage tubulaire de $D$, ou, ce qui revient au même, comme le bord du
disque $\Delta(D)$ complémentaire de ce voisinage ; et la carte duale s'obtient
en recollant les disques $\Delta(D)$ d'une façon naturelle, les « arêtes
libres » formant le bord. En gardant la trace, à l'intérieur de $\Delta(D)$, du
système de pseudo-droites, on obtient un disque $\Delta'(D)$ où les transitions
apparaissent comme des segments qui vont du bord vers l'intérieur ; mais les
identifications qui y mènent créent des points doubles artificiels sur le bord :
des pseudo-droites qui ne se rencontrent pas dans $\Delta(D)$ se rencontrent sur
le bord de $\Delta'(D)$.

Si l'on recolle les $\Delta(D)$ \emph{sans} faire ces identifications, on obtient
une surface à bord, dont les points du bord sont d'indice $2$ ou $3$ et dont
chaque sommet est sur le bord ; en bouchant ses trous par des disques
$\Delta_x$, on retrouve une surface compacte sans bord, et l'on retrouve la
surface du début en contractant ces disques en des points. Les trous de la
surface à bord correspondent canoniquement aux sommets de la surface sans
bord.\footnote{Page 27. La page 29 est une page de dessins sans texte : cinq
disques traversés de pseudo-droites, accolés en chaîne autour d'un point où
leurs bords se rejoignent — le recollement des $\Delta(D)$ que décrivent les
pages 25 et 27.}

\subsection*{Vers une position supercritique (25 décembre)}

Reste à définir le polygone d'une position supercritique $\Delta \in \Sigma$,
comme l'ensemble de ses générisations muni d'une structure polygonale
convenable. On a vu (page 14) que ces générisations correspondent aux sommets
de $\widetilde{\Delta}$ ; et pour un sommet $s$ de $\Delta$, les deux points de
$\widetilde{\Delta}$ au-dessus de $s$ correspondent aux deux orientations de $X$
en $s$, c'est-à-dire aux deux \emph{rives} de $\Delta \smallsetminus \lbrace s
\rbrace$. Une spécialisation $D \to \Delta$ choisit l'une des rives, celle
qu'elle balaie, et l'on trouve un isomorphisme combinatoire canonique entre le
segment qui correspond à la rive balayée et un segment $\mathrm{I}(\Delta, s)$
découpé sur $\widetilde{\Delta}$ par l'antipodisme.\footnote{Page 28, dont la
moitié inférieure est couverte d'un grand encadré barré et de lignes biffées ;
la transcription n'en donne que ce qui se lit, et la définition de
$\mathrm{I}(D', s)$ n'est pas lisible. La page écrit $D'$ pour la position
supercritique.} La construction s'arrête là ; la question de la spécialisation
vers une position supercritique revient aux pages 40 et 54.

\pagerange{30}{33}
\section*{Les cartes en général (pages 30 à 33)}

Les pages 30 à 33 isolent, hors de toute pseudo-droite, la construction qui
fait d'un graphe muni de polygones locaux une surface. C'est le formalisme des
cartes combinatoires, et il est écrit ici avec une netteté qui permet de le
redonner tel quel.

\subsection*{La donnée, et la surface}

Soit $\Gamma$ un graphe : un ensemble $S$ de sommets, un ensemble $\vec{A}$
d'arêtes orientées muni de l'involution sans point fixe $\vec{a} \mapsto
-\vec{a}$ et de l'application origine. Pour $s \in S$, soit $\vec{A}_s$
l'ensemble des arêtes d'origine $s$. Une \emph{semi-circularisation} de $\Gamma$
est la donnée, pour tout $s$, d'une structure polygonale sur $\vec{A}_s$ — un
ordre cyclique défini au sens près. Soit $\omega_s = \omega(\vec{A}_s)$
l'ensemble de ses deux orientations. On se donne enfin, pour toute arête
$\vec{a} : s \to t$, une bijection de transport
\[
\varphi_{\vec{a}} : \omega_s \longrightarrow \omega_t , \qquad
\varphi_{-\vec{a}} = \varphi_{\vec{a}}^{-1} .
\]
On suppose, pour simplifier, que les sommets sont d'ordre au moins
$3$.\footnote{Pour les sommets d'ordre $1$ ou $2$, il faut se donner de plus un
polygone combinatoire contenant $\vec{A}_s$ et une orientation ; en marge : « cas
$\nu_s = 1, 2$ n'est pas exclu ».}

En termes d'aujourd'hui, c'est un \emph{système de rotation signé} : l'ordre
cyclique en chaque sommet est donné au sens près, et les $\varphi_{\vec{a}}$
disent, arête par arête, si les deux sens se raccordent ou se
croisent.\footnote{La notion est classique en théorie topologique des graphes
(Heffter, puis Edmonds en 1960, pour le cas orientable ; Ringel, puis Stahl en
1978, pour les surfaces non orientables, sous le nom de schéma de plongement).
La version du dossier est intrinsèque — elle ne choisit pas d'orientation
locale. Le formalisme des drapeaux qui suit est celui des « cartes » de
l'\emph{Esquisse d'un programme}, que Grothendieck rédige en janvier 1984, dans
les mêmes semaines ; aucune des deux ne renvoie à l'autre.}

La surface se construit ainsi. Pour chaque sommet $s$, soit
$\operatorname{Dis}_s$ un disque polygonal réalisant le polygone \emph{dual} de
$\vec{A}_s$ : ses côtés $\alpha_{\vec{a}}$ correspondent aux arêtes $\vec{a}$
issues de $s$, et il se décompose en $2\nu_s$ drapeaux, un par repère de
$\vec{A}_s$. On recolle alors le côté $\alpha_{\vec{a}}$ de
$\operatorname{Dis}_s$ au côté $\alpha_{-\vec{a}}$ de $\operatorname{Dis}_t$ par
l'homéomorphisme $\psi_{\vec{a}}$ qui \emph{renverse} les orientations induites
par $\omega \in \omega_s$ et par $\varphi_{\vec{a}}(\omega) \in \omega_t$. On
obtient une surface topologique compacte sans bord $\mathcal{X}$, dans laquelle
la réalisation $|\Gamma|$ est canoniquement plongée, les disques
$\operatorname{Dis}_s$ étant les faces de la carte duale $\mathcal{C}^{*}$.

\subsection*{Les repères et les trois involutions}

Soit $\operatorname{Rep}$ l'ensemble des \emph{repères}, les couples
$(\vec{a}, \omega)$ avec $\vec{a} \in \vec{A}_s$ et $\omega \in \omega_s$. Pour
$\omega \in \omega_s$, soit $\rho_{\omega}$ la permutation circulaire de
$\vec{A}_s$ qu'elle définit. La page 31 pose
\[
\sigma_0(\vec{a}, \omega) = \bigl(-\vec{a}, -\varphi_{\vec{a}}(\omega)\bigr),
\qquad
\sigma_1(\vec{a}, \omega) = \bigl(\rho_{\omega}(\vec{a}), -\omega\bigr),
\qquad
\sigma_2(\vec{a}, \omega) = (\vec{a}, -\omega).
\]
Ce sont trois involutions — pour $\sigma_1$, parce que $\rho_{-\omega} =
\rho_{\omega}^{-1}$ — et $\sigma_0$, $\sigma_2$ commutent. Leurs composés sont
les rotations de la carte $\mathcal{C}$ :
\[
\rho_{\mathrm{som}} = \sigma_2\sigma_1 : (\vec{a}, \omega) \longmapsto
(\rho_{\omega}(\vec{a}), \omega), \qquad
\rho_{\mathrm{face}} = \sigma_1\sigma_0 : (\vec{a}, \omega) \longmapsto \bigl(
\rho_{\omega_t}^{-1}(-\vec{a}), \omega_t \bigr), \ \omega_t =
\varphi_{\vec{a}}(\omega),
\]
la rotation autour du sommet et le parcours du bord d'une face : les faces
orientées de $\mathcal{C}$ sont les orbites de $\rho_{\mathrm{face}}$, comme le
dit la page 31 en premier lieu. Le troisième composé,
\[
\sigma = \sigma_2\sigma_0 = \sigma_0\sigma_2 : (\vec{a}, \omega) \longmapsto
\bigl(-\vec{a}, \varphi_{\vec{a}}(\omega)\bigr),
\]
est l'involution qui retourne une arête en gardant l'orientation
locale.\footnote{La dernière ligne de la page 31 écrit « $\rho_f = \sigma_2\sigma_0
\overset{\mathrm{def}}{=} \sigma$, $\sigma(\vec{a}, \omega_s) = (-\vec{a},
\varphi_{\vec{a}}\,\omega_s)$ », le $\rho_f$ récrit sur un « $\rho_{\mathrm{face}}$ »
biffé. La formule de $\sigma$ est juste ; le nom $\rho_f$ ne l'est pas, puisque
le parcours d'une face défini au début de la même page est
$\sigma_1\sigma_0$, comme la ligne précédente l'écrit, et comme le reprend la
page 156 (« $\rho_f = \sigma_1\sigma_0$, $\rho_s\rho_f = \sigma$ »). Nous
suivons la page 156.} C'est, lettre pour lettre, la présentation d'une carte par
son ensemble de drapeaux et le \emph{groupe cartographique} $\langle \sigma_0,
\sigma_1, \sigma_2 \mid \sigma_i^2 = (\sigma_0\sigma_2)^2 = 1 \rangle$ : $\sigma_i$
change la composante de dimension $i$ du drapeau.\footnote{En identifiant
$\operatorname{Rep}(\mathcal{C})$ à $\coprod_s \operatorname{Rep}(\vec{A}_s)$, la
page note que $\sigma_1$ et $\sigma_2$ sont les opérations habituelles
$\widetilde{\sigma}_0^{s}$, $\widetilde{\sigma}_1^{s}$ sur les repères du polygone
$\vec{A}_s$, et que $\rho_{\mathrm{som}}$ est la rotation $\widetilde{\rho}^{s}$
de ce polygone. La présentation d'une carte par trois involutions est due à
Tutte (1973) et, sous forme de graphes à arêtes coloriées, à Lins (1982) ; le
groupe cartographique est le nom que lui donne l'\emph{Esquisse}.}

\subsection*{La surface à bord, et le cas des pseudo-droites}

La subdivision barycentrique de $\mathcal{C}$, qui est aussi celle de
$\mathcal{C}^{*}$, a pour faces les drapeaux. Une seconde construction donne la
surface privée de petits disques autour des centres des faces de $\mathcal{C}$ :
on « enlève les coins » de chaque $\operatorname{Dis}_s$, ce qui donne un
polygone $\operatorname{Dis}'_s$ d'ordre $2\nu_s$, dont les côtés alternent
entre côtés de recollement, qui correspondent aux arêtes issues de $s$, et côtés
libres, qui correspondent aux coins. En ne recollant que les côtés de
recollement, on obtient une carte à bord dont les faces sont les
$\operatorname{Dis}'_s$ et dont les composantes du bord correspondent aux faces
de $\mathcal{C}$, c'est-à-dire aux orbites du groupe diédral $\langle \sigma_0,
\sigma_1 \rangle$, qui contient $\langle \rho_{\mathrm{face}} \rangle$ comme
sous-groupe d'indice $1$ ou $2$. Chaque composante du bord est un polygone,
canoniquement isomorphe au dual du polygone de la face correspondante.

Dans le cas qui intéresse le dossier — la carte de toutes les positions
relatives d'un système de pseudo-droites —, les faces sont d'ordre divisant
$4$, et parmi les quatre côtés de chaque composante du bord, qui correspond à un
coin de la carte, l'un est privilégié. Mais ce qui le distingue ne dépend pas
seulement du sommet, et dépend du coin : on ne peut donc pas espérer, en
général, réaliser $\mathcal{C}$ comme la réunion de deux graphes $\gamma$,
$\gamma^{*}$ associés à deux cartes duales l'une de l'autre sur la
surface.\footnote{Page 33, dont la fin se lit mal. Les figures des pages 30 à 32
superposent en vert foncé la carte et en vert clair sa duale, les drapeaux
hachurés, et marquent $r$, $\sigma_0 r$, $\sigma_1 r$, $\sigma_2 r$ autour d'un
sommet ; la page 32 appelle « bleues » des arêtes tracées en vert clair.}

\pagerange{33}{42}
\section*{Le disque d'une position, ses arcs, et le recollement (pages 33 à 42)}

\subsection*{Le disque $\Delta_L$ et son système de cordes}

Soit $L$ une position relative non triviale. Une bande de Möbius étroite $B_L$
autour de $L$ a pour bord le revêtement des rives $\widetilde{L}$, et le
complémentaire de son intérieur est un disque $\Delta_L$. Toute $D_i$ coupe
$\partial\Delta_L = \widetilde{L}$ en exactement deux points, et sa trace
$d_i = D_i \cap \Delta_L$ est un segment. Le système de cordes $(d_i)$ a les
propriétés suivantes.
\begin{itemize}
\item[a)] Deux cordes se coupent en au plus un point, intérieur à $\Delta_L$, et
s'y croisent.
\item[b)] La relation « $d_i$ et $d_j$ ne se croisent pas dans $\Delta_L$ » — le
\emph{parallélisme} — est une relation d'équivalence ; une classe de
$\nu \geqslant 2$ cordes a ses $2\nu$ extrémités réparties en deux segments
antipodiques $S'$, $S''$ de $\nu$ points consécutifs de
$\partial\Delta_L$.\footnote{La page écrit la relation « $d_i = d_j$ ou
$\partial d_i \cap \partial d_j = \emptyset$ », qui, prise à la lettre, rendrait
parallèles presque toutes les cordes. Nous lisons ce que la suite demande :
deux cordes sont parallèles quand $D_i \cap D_j$ tombe dans la bande, c'est-à-dire
sur un sommet de $\Sigma$ situé sur $L$ ; la relation est alors transitive
parce que $D_j$ ne rencontre $L$ qu'une fois, et ses classes sont les faisceaux
de pseudo-droites qui passent par un même sommet de $L$. L'alinéa b) est d'une
écriture très rapide, avec des crochets de sa main qui ne se referment pas.}
\end{itemize}
Ces conditions combinatoires, nécessaires pour que le système de cordes soit
celui d'une position relative, sont données comme suffisantes — sans
démonstration.\footnote{« Ces conditions combinatoires nécessaires (\ldots) sont
aussi suffisantes » (page 34). Nous n'avons pas trouvé, dans le dossier,
d'argument ; la réciproque demande de recoller une bande de Möbius au bord du
disque et de vérifier que deux cordes parallèles s'y rencontrent exactement une
fois, ce que la condition b) semble faite pour garantir.} En choisissant un
représentant dans chaque classe de parallélisme, on obtient un système de cordes
strict, où deux cordes quelconques se croisent — l'image d'un arrangement de
pseudo-droites dans un disque, celle qu'étudie le dossier 80.

\subsection*{Les arcs de transition}

Sur le bord $\widetilde{L}$, plusieurs familles d'arcs.
\begin{itemize}
\item Les \emph{arcs de générisation} : les deux arcs interceptés par $S'$ et
$S''$ pour chaque classe de parallélisme ; à l'antipodisme près, ils sont en
bijection avec $L \cap S$.
\item Les \emph{arcs de spécialisation ordinaires} : on considère les ensembles
$\sigma$ de sommets du bord qui forment un arc combinatoire, avec
$\operatorname{card}\sigma \geqslant 2$, tels que les cordes issues des points
de $\sigma$ passent par un même sommet $t$ de $\Sigma$ et épuisent les cordes
qui passent par $t$, et dont les extrémités ne sont pas dans un arc de
générisation. Pour deux points voisins $s$, $s'$ de $\sigma$ il existe alors un
unique triangle de $\Sigma$ de sommets $t$, $s$, $s'$ ; leur réunion est un
triangle de spécialisation, et sa base l'arc de spécialisation.
\item Les \emph{arcs de spécialisation exceptionnels}, qui correspondent aux
couloirs de $\Sigma \cup \lbrace L \rbrace$ adjacents à $L$ dont le sommet est
un sommet de $\Sigma$ : ce sont des arcs ouverts de $\widetilde{L}$ délimités par
une classe de parallélisme. Ils ne rencontrent aucun arc de spécialisation
ordinaire ni aucun arc de générisation.
\end{itemize}
$L$ a deux arcs exceptionnels si et seulement si elle est la « bissectrice »
d'un couloir ; elle n'a alors pas d'autre arc de spécialisation et exactement
deux arcs de générisation, et ces quatre arcs sont les arêtes d'une structure de
polygone sur $\widetilde{L}$. Dans tous les cas, deux \emph{arcs de transition}
distincts (de spécialisation ou de générisation) ne se rencontrent que si l'un
est un arc de spécialisation exceptionnel et l'autre un arc de générisation, et
ils sont alors adjacents.

Les \emph{sommets de transition}, extrémités des arcs de transition, découpent
$\widetilde{L}$ en arcs de trois types : de générisation, de spécialisation, et
les autres, les \emph{arcs caractéristiques} ou \emph{arcs d'inertie}.
Quitte à rétrécir un peu les arcs exceptionnels — à les arrêter aux premiers
sommets qui suivent leurs extrémités —, les arcs de transition deviennent deux à
deux disjoints, et de deux arcs caractéristiques adjacents un et un seul est un
arc d'inertie.\footnote{Pages 36-37. La page 37 appelle aussi cet arc « arc de
neutralité », lecture probable. La fin de la page 36, rapide et raturée, se lit
en partie seulement.}

Les arcs de générisation servent à ce que leur nom dit : la générisation
consiste à « écarter les brins » au-dessus de l'arc. Dans $X$, c'est un
\emph{décollement} du sommet $s$ de $\Sigma$ dans la direction de l'arc, qui
« pousse » sur $L$ ; le polygone $\widetilde{L}$, à isomorphisme combinatoire
canonique près, n'en est pas changé.\footnote{Au-dessus de « décollement », un
mot lu « antibalayage », d'une lecture douteuse. Figures de la page 37 : trois
disques reliés par des flèches, les extrémités de deux cordes $D$, $D'$ écartées
en haut et croisées en bas.}

\subsection*{Le recollement, première version}

Si $(L, a)$ et $(L', a')$ se correspondent — $a$ arc de générisation de $L$,
$a'$ arc de spécialisation de $L'$ —, on attache $\Delta_L$ à $\Delta_{L'}$ en
identifiant $a$ à $a'$, par l'application qui envoie chaque extrémité de l'un
sur l'extrémité de l'autre qui appartient à la même pseudo-droite, donc en
intervertissant les extrémités. Mais, « réflexion faite », ce recollement
correspond, sur les orientations $\omega(\Delta_L) \simeq \omega(L)$ et
$\omega(\Delta_{L'}) \simeq \omega(L')$, à l'isomorphisme évident, et ce n'est
pas le bon : il faut prendre celui qui fixe $s$ et $t$ et échange $D$ et $D'$,
la symétrie par rapport à l'axe $(s, t) = a$. Le transport des orientations est
alors l'\emph{opposé} de l'isomorphisme tautologique.\footnote{Pages 37-38. Ce
choix sera repris et discuté entièrement aux pages 64 à 70, où Grothendieck
conclut qu'il s'était trompé deux fois.}

Avec cette identification, toutes les faces de la carte obtenue sont d'ordre
$4$. Leurs coins, qui viennent des arcs d'inertie, ont des côtés de trois types
— spécialisation, générisation, mixte — et, en faisant le tour d'une face à
partir du côté de spécialisation, on rencontre dans l'ordre les transitions
\begin{gather*}
\text{spécialisation} \to \text{spécialisation}, \qquad \text{générisation} \to
\text{spécialisation}, \\
\text{générisation} - \text{générisation}, \qquad
\text{spécialisation} \to \text{générisation},
\end{gather*}
soit une arête du type spécialisation, une du type générisation et deux du type
mixte. Les faces correspondent bijectivement aux couples $(L, a)$ où $a$ est un
arc d'inertie du type spécialisation — ou du type générisation, les deux
descriptions se correspondant.\footnote{La page 175, la dernière du dossier,
dessine exactement un tel carré : « $1 \xrightarrow{\text{spéc}} 2
\xrightarrow{\text{gén}} 3 \xrightarrow{\text{gén}} 4 \xrightarrow{\text{spéc}} 5
= 1$ ». La seconde description, « peut-être plus jolie », revient à chercher les
paires $(s, t)$ de sommets de $\Sigma$, les positions qui y passent, et un arc
de neutralité au-dessus de $(s, t)$ parmi quatre candidats ; elle reste « à
expliciter ».}

\subsection*{Les trous, et les positions supercritiques}

La construction n'est pas complète : on a laissé de côté les positions
supercritiques et les arcs qui y mènent. Après avoir bouché les « trous
quadrangulaires » — circuits de quatre arcs d'inertie —, il reste des trous à
boucher par les disques des positions supercritiques, qu'il faut recoller le
long des arcs de spécialisation exceptionnels, restés libres. La page 40 en
donne le cadre : un graphe $\Gamma = (S, \vec{A}, \sigma)$ dont certaines arêtes
sont \emph{libres}, laissées fixes par l'involution $\sigma$, « arêtes orientées
sans extrémité » ; en recollant les $\operatorname{Dis}_s$ le long des arêtes
ordinaires seulement, on obtient toujours une surface à bord. Il faut donc
étudier les circuits formés d'arcs de spécialisation exceptionnels, et voir
comment les boucher.\footnote{La page 40 écrit « $\Gamma = (S, \vec{A}, \sigma,
\sigma)$ », les deux dernières lettres semblables ; nous gardons une seule
involution.}

Pour $L = D_i \in \Sigma$, on définit encore $\Delta_L$ en voyant $D_i$ comme
position relative par rapport à $\Sigma \smallsetminus \lbrace D_i \rbrace$ : un
disque avec $n - 1$ cordes, et un polygone d'ordre $2(n-1)$. La spécialisation
d'une position $L$ vers $D$ le long d'un couloir revient, sur $\Delta_L$, à
« tirer » l'arc de spécialisation par-dessus $D$ ; la configuration obtenue est
canoniquement celle de $(D, \Delta_D)$, et l'on recolle $\Delta_L$ à $\Delta_D$
par cet isomorphisme, dont le transport des orientations est à nouveau l'opposé
de l'isomorphisme tautologique. L'ennui est que les arcs de générisation de $D$
empiètent les uns sur les autres : chaque sommet de $D$ en détermine deux.

Remède : « multiplier » chaque $D \in \Sigma$ en autant d'exemplaires qu'elle a
de sommets. Pour un couple $(D, s)$ il y a alors exactement deux arcs de
générisation exceptionnels, qui n'empiètent pas : si $s$ est un sommet double,
ce sont les deux côtés d'un bigone ; sinon, avec les arcs ordinaires, ils
forment un polygone. Dans les deux cas il n'y a pas d'arc d'inertie — « ce qui
est un triomphe ! » — et l'on complète la surface en une surface compacte en
recollant des bigones ou des carrés.\footnote{Page 42. La phrase suivante,
inachevée en bas de page, propose de remplacer dans le graphe chaque sommet
exceptionnel par l'ensemble des diagonales du polygone local.}

\pagerange{43}{46}
\section*{Le déploiement en axiomes (pages 43 à 54, 31 décembre 1983 au 2 janvier 1984)}

Le 31 décembre, une nouvelle rédaction, titrée « Déploiement d'un système de
pseudo-droites », renverse la démarche : au lieu de construire la surface, on
en écrit les propriétés, de sorte qu'on puisse la reconnaître et, à la fin,
reconstituer $\Sigma$ à partir d'elle.\footnote{Le titre porte, dans la
transcription, un numéro de section de LaTeX ; il n'est pas de sa main, et la
rédaction n'a pas de numéro de page. Les pages 43 à 45 sont denses d'ajouts
interlinéaires, et plusieurs axiomes ne se lisent qu'en partie.}

\subsection*{Les axiomes}

Un \emph{déploiement} de $\Sigma$ est une surface compacte $\mathcal{X} =
\mathcal{X}(\Sigma)$ munie de deux graphes $K$ et $L$ qui la découpent
cellulairement, avec les propriétés suivantes.
\begin{itemize}
\item[a)] Tous les sommets de $K$ sont d'ordre $4$ ; ceux de $L$ sont d'ordre
pair.
\item[b)] Toute arête de $K$ contient au plus un sommet de $L$.
\item[c)] La réunion $K_0$ des adhérences des arêtes de $K$ qui ne contiennent
aucun sommet de $L$ — les \emph{arêtes exceptionnelles} — est une réunion finie
de cercles. Il revient au même de dire que les voisins d'une arête
exceptionnelle sont exceptionnels et que de tout sommet de $K$ part au moins une
arête ordinaire.
\item[d)] Soit $K_1$ la réunion des arêtes ordinaires, de sorte que $K = K_0
\cup K_1$. On se donne, pour chaque arête ordinaire, une \emph{rive
privilégiée} : une section du revêtement des bords de l'immersion du déployé
$K'_1$ dans $\mathcal{X}$.\footnote{En marge : « cette donnée résulte en fait
des autres, (\ldots) ; cf.\ f) plus bas ».}
\item[e)] Un sommet exceptionnel — incident à une arête exceptionnelle — porte
une multiplicité, lue sur les arêtes ordinaires qui en partent.\footnote{Alinéa
lu en partie seulement.}
\item[f)] En tout point $x \in K \cap L$, les branches de $K$ et de $L$ se
croisent, et trois cas se présentent : $x$ n'est pas un sommet de $L$ (il n'est
alors pas un sommet de $K$) ; $x$ est un sommet de $L$ mais pas de $K$ ; $x$ est
un sommet des deux. Un sommet exceptionnel de $K$ n'est pas sur $L$ (f').
\item[g)] Pour un sommet ordinaire $s$ de $K$, le nombre de points de $L$ entre
$s$ et l'unique sommet de $L$ sur chacune des quatre arêtes issues de $s$ ne
dépend pas de l'arête ; on le note $\nu(s)$, et l'on pose $\varepsilon(s) = 1$
si $s \in L$, $\varepsilon(s) = 0$ sinon.\footnote{La définition de
$\varepsilon(s)$ est biffée à la page 43 ; elle revient telle quelle à la page
170 (alinéa g), et nous la gardons.}
\item[h)] Pour toute face $F$ de $K$, le nombre de points de $L$ sur $\partial F$,
comptés avec leur multiplicité réduite, est $2n$ ; pour une face ordinaire :
\[
\sum_{s \text{ sommet de } F} \bigl( 2\nu(s) + \varepsilon(s) \bigr) \;+\;
\sum_{a \text{ arête de } F} \nu(a) \;=\; 2n .
\]
\end{itemize}
Le sens de ces axiomes : une face de $K$ est le disque d'une position, son bord
porte les $2n$ points du polygone $\operatorname{Pol}(D)$, et $L$ est la trace
des pseudo-droites de $\Sigma$.

\subsection*{L'indexation}

Les cellules de $(\mathcal{X}, K)$ sont indexées par la géométrie de
$\Sigma$ :
\begin{itemize}
\item les faces, par les positions relatives non supercritiques ; on note $F_D$
la face de $D$ ;
\item les arêtes, par les couples $(D, D')$ où $D \to D'$ est une spécialisation
immédiate, l'arête étant ordinaire si et seulement si c'est une vraie
spécialisation, et sa rive privilégiée celle de $F_D$ ;
\item les sommets ordinaires, par les quadruplets $(D, t, t', \beta)$ où $D$ est
une position non supercritique, $t$ et $t'$ deux triangles de spécialisation
pour $D$ et $\beta$ l'arc de $\widetilde{D}$ compris entre leurs arcs : les
quatre faces autour du sommet sont $D$, ses spécialisations $D'$, $D''$ à
travers $t$ et à travers $t'$, et $D'''$ à travers les deux à la
fois ;\footnote{Les lettres des triangles varient sur la page ($t'$, $t''$ dans
le quadruplet, $t$, $t'$ ensuite), et la page semble écarter, dans un ajout,
le cas où $D$ n'est pas stable.}
\item les sommets exceptionnels, par les drapeaux $(s, a)$ de type $(0, 1)$ de
$\Sigma$ : un sommet $s$ et une arête $a$ d'une $D_i$ en $s$.
\end{itemize}
« Ces indications suffisent à construire le tout cellulaire $(\mathcal{X}, K)$,
qui est une géométrie d'incidence. »

\textbf{Corollaire (page 44).} \emph{Les composantes connexes de $K_0$ sont
indexées par $I$, et la composante $K(i)$ est combinatoirement isomorphe à
$\widetilde{D}_i$, l'isomorphisme dépendant du choix d'une orientation de
$D_i$.}

\subsection*{L'application structurale}

(i) On a une application $\varphi : \mathcal{X} \to X$, « l'application
structurale », qui au voisinage des arêtes de $K$ a la structure d'un
« folding » (un pli) et au voisinage des sommets de $K$ celle d'un « double
folding », et telle que
\[
K \cap L = (\varphi|_K)^{-1}(X_1) .
\]
La page écrit aussi $L = \varphi^{-1}(X_1)$ ; ce n'est vrai qu'en dehors de
$K_0$, puisque chaque cercle exceptionnel $K(i)$ revêt $D_i$ par le corollaire
précédent. L'énoncé juste est
\[
\varphi^{-1}(X_1) = L \cup K_0 ,
\]
et c'est celui que donne la page 159, le lendemain.\footnote{Page 45, (i) 4) :
« $L = \varphi^{-1}(\bigcup D_i)$ », la réunion notée $X_1$. La page 159, datée
du 1er janvier 1984 et barrée d'une croix verte, écrit
« $\varphi^{-1}(X_1) = L \cup K_0$ ». La propriété 1), le comportement de
$\varphi$ sur chaque face, est illisible.}

\textbf{Corollaire (page 45).} \emph{Un point $s \in K \cap L$ a pour image un
sommet de $\Sigma$ si et seulement si c'est un \emph{point critique} de $K$.}

(j) Soit $F$ une face. Aux points critiques $s$ de $K$ sur $\partial F$ dont la
rive privilégiée est extérieure à $F$, corrigeons le contour de $F$ par un petit
renflement du côté privilégié, et soit $F' \supset F$ la face ainsi corrigée.
Alors la position relative de $F$ est $\varphi(\partial F')$, et la situation
affine correspondante — le disque $\Delta_D$ et ses cordes — est canoniquement
isomorphe à $(F', F' \cap L)$.

k) à n) Une arête exceptionnelle est stable ou instable selon que ses deux
positions incidentes sont stables ou non ; en un sommet exceptionnel, la rive
privilégiée de la branche ordinaire est du côté de l'arête instable. Une loi de
symétrie règle la structure autour d'une arête exceptionnelle, et l'indice des
arêtes exceptionnelles se déduit des autres données numériques.\footnote{Page
46 : l'indice d'une arête exceptionnelle stable est « $n - (\nu_{s} - 1) -
(\nu_t - 1)$ », celui d'une instable « $n - \nu$ ». Les pages 55, 93 et 110
comptent la même chose, semble-t-il, comme $(n - 1) - (\nu - 1) - (\nu' - 1)$,
une unité de moins ; nous ne savons pas laquelle des deux quantités la page 46
appelle indice, et nous ne tranchons pas.} Les sections de $L$ sur les faces
sont indexées par $\widetilde{I}$ (m), et l'on doit pouvoir reconstituer
$\Sigma$ à partir de $K$ et des données numériques $\varepsilon(s) \in \lbrace
0, 1 \rbrace$, $\nu(s) \in \mathbf{N}$, $\nu(a)$ (n). Enfin il y a aussi une loi
de symétrie autour d'une arête ordinaire $a$ : si $F'$, $F''$ sont les deux
faces incidentes et $T$ un voisinage tubulaire de $K \cup L$, une involution
$\sigma_a$ de $T_a = (F' \cup F'') \cap T$ qui fixe $a$ point par point et
respecte la structure cellulaire induite par $K \cup L$ ; pour deux arêtes
incidentes au même sommet, $\sigma_a$ et $\sigma_b$ commutent sur $T_a \cap T_b$.

\pagerange{47}{54}
\subsection*{La structure de $K$ : sommets rouges et noirs, signes (pages 47 à 54)}

On ajoute à $K$ des sommets d'ordre $2$ ; l'ensemble des sommets devient $S_K =
S'_K \amalg S''_K$, sommets « rouges » et « noirs » (ou orangés), comme sur les
figures. Deux rédactions du placement de ces sommets, toutes deux du 2 janvier,
se suivent aux pages 48 et 50 ; la seconde s'énonce ainsi : pour toute face $F$,
$\partial F \cap S'_K$ est un polygone d'ordre $2n$, et chacune de ses arêtes
porte un point de $S''_K$ et un seul ; les points rouges sur $L$ sont les points
de $K \cap L$ autres que les points critiques d'ordre pair ; toute arête de $K$
rencontre $S'_K$ ; et près d'un point critique $s$, le segment entre les deux
points de $L$ les plus proches contient $\nu(s)$ points rouges.\footnote{Page 50,
datée « 2.1.84. » ; la page 48, datée « 2.1.84 », en est un premier état, et
toutes deux sont traversées de deux longs traits obliques qui ne masquent pas
le texte.}

\textbf{Les signes.} Chaque arête $a$ d'une face $F$ porte un signe
$\varepsilon_F(a) = \pm 1$, et l'on a les règles : $\varepsilon_{F'}(a) =
-\varepsilon_F(a)$ si $F$ et $F'$ sont les deux faces incidentes à $a$ ;
$\varepsilon_F(a) = \varepsilon_F(b)$ si $a$ et $b$ sont adjacentes sur $F$ en un
sommet d'ordre $2$ ; et $\varepsilon_{F_1}(a_1) = \varepsilon_{F_2}(a_2)$ si
$F_1$, $F_2$ sont adjacentes le long d'une arête $b$ d'extrémité $s$ d'ordre $4$,
$a_i$ étant l'autre arête de $F_i$ en $s$.\footnote{Page 47, datée « 2.2.84 »
(vérifié sur le fac-similé par Michel Hua ; la première lecture était
« 7.1.1984 ») : le feuillet est donc postérieur d'un mois à ceux qui
l'entourent. La troisième règle est écrite $\varepsilon_{a_1}(F_1) =
\varepsilon_{a_2}(F_2)$ ; nous rétablissons l'ordre des indices des deux
premières.}

\textbf{Les bijections de symétrie.} Toutes les faces ayant un polygone
$\operatorname{Pol}(F)$ du même ordre $2n$, une arête $a$ commune à $F$ et $F'$
définit une \emph{bijection de symétrie}
\[
\sigma_a^F : \operatorname{Pol}(F) \xrightarrow{\ \sim\ } \operatorname{Pol}(F'),
\qquad \sigma_a^F \sigma_a^{F'} = \mathrm{id},
\]
caractérisée par le fait qu'elle fixe les deux extrémités de $a$ ; elle fixe
alors tout ce qui est sur la « grande arête » qui contient $a$, entre les
sommets d'ordre $4$ les plus proches. « J'aimerais comprendre la relation entre
les $\sigma_a^F$ et les signes $\varepsilon_F(a)$ » (page 49) : c'est la
question que les pages suivantes attaquent, sans la résoudre
entièrement.\footnote{C'est aussi la transcription discrète des transports
$\sigma_{a_i}$ de la page 55, qui lisent chacun comme une réflexion
$t \mapsto l_i - t$ de $\mathbb{Z}/2n\mathbb{Z}$.}

\textbf{Les points critiques.} Soit $C_K \subset S_K$ l'ensemble des points
critiques : des sommets d'ordre $2$, un exactement sur chaque grande arête non
exceptionnelle, munis d'un poids $\nu(s) \geqslant 1$ dont la parité dit la
couleur.\footnote{Page 49 : « $\nu(s) \geqslant 1$ », le $1$ surchargé,
peut-être un $0$ ; le poids d'un point non critique est $1$ s'il est rouge, $0$
s'il est noir. L'ajout sur la parité est souligné et se lit mal.} On a : si $s$
est critique sur $F$ et $s'$ son antipodique sur $\partial F$, le signe vaut
$+1$ si et seulement si $s'$ est aussi critique ; il y a au moins trois points
critiques sur tout $\partial F$, et au moins un de signe $+1$ ; de chacun des
deux arcs ouverts découpés par $s$ et son antipodique, l'un au moins contient un
point critique. La page 51 compare les points critiques de $F$ et de $F' =
\sigma_a(F)$ le long d'une grande arête $a$ : une bijection entre les points
critiques de deux arcs qui se correspondent, qui conserve les signes sauf au
point $s$ de $a$, où il change ; et un nombre $N'$ de sommets d'un certain arc
de $F'$ compris entre $N - 1$ et $N + 1$.\footnote{Le passage est chargé
d'ajouts et de ratures et se lit mal. Le \emph{corollaire de la page 51} dit
qu'un certain arc $\beta^{*} = \widehat{u^{*} s\, v^{*}}$, entre les
antipodiques des deux points critiques les plus proches de $s$, ne contient pas
de point critique d'un signe donné, la fin de l'énoncé étant illisible ;
« En résumé » distingue ensuite trois cas, $x = s$, $x = t$ et $x$ intérieur à
$\beta$.}

\textbf{Proposition (page 53).} \emph{Soit $a$ une grande arête de la face $F$,
et $F' = \sigma_a(F)$ la face adjacente le long de $a$. Soit $A$ la réunion des
grandes arêtes $\alpha$ de $F$ telles que $\sigma_a(\alpha)$ soit une grande
arête de $F'$.}
\begin{itemize}
\item[a)] \emph{$A$ est un arc de $\partial F$. Soit $B$ l'arc complémentaire ;
les nombres de grandes arêtes $n_a(F)$ et $n_a(F')$ qu'il porte, sur $F$ et sur
$F'$, sont compris entre $3$ et $5$, et les possibilités sont, à l'échange de $F$
et $F'$ près, $(3, 3)$, $(3, 4)$ et $(3, 5)$. Dans les cas inégaux, c'est la
face spéciale qui porte $3$ arêtes, ce qui suffit à la déterminer.}
\item[b)] \emph{Pour une petite arête $\alpha \neq a$ de $F$ contenue dans $A$,
$\varepsilon_{F'}(\alpha) = \varepsilon_F(\alpha)$ ; pour $\alpha = a$,
$\varepsilon_{F'}(a) = -\varepsilon_F(a)$.}
\item[c)] \emph{Les distributions possibles des grandes arêtes sur $B$ et de
leurs signes sont les trois configurations $1^{\circ}$, $2^{\circ}$,
$3^{\circ}$ dessinées sur la page.}
\end{itemize}
Les trois cas sont numérotés sur la page sans être rattachés aux trois couples ;
d'après la page 54, le cas $2^{\circ}$ est le cas « symétrique », donc $(3, 3)$,
et les cas $1^{\circ}$ et $3^{\circ}$ sont les deux cas inégaux.\footnote{Les
couples, surchargés, sont lus $(3, 3)$, $(3, 4)$, $(3, 5)$, les seconds chiffres
d'après « compris entre $3$ et $5$ » ; un « $(2, 2)$ » est biffé dans un crochet.
Les figures des pages 53 et 56 montrent des arcs verts d'extrémités $u$, $v$,
avec les signes $+1$ au sommet et $-1$ entre les sommets foncés. La page 52 est
une page de dessins sans texte : le bord $\partial F$ et ses sommets critiques.}

\textbf{Corollaire (page 54).} \emph{La distribution des signes sur $F$ détermine
celle sur $F'$.} C'est clair dans les cas $1^{\circ}$ et $3^{\circ}$, où l'on
sait laquelle des deux faces est la spécialisation de l'autre ; dans le cas
symétrique, les deux grandes arêtes extrêmes ont les mêmes signes sur $F$ et sur
$F'$, et celles qui vont par paires ont des signes opposés.

\textbf{Second corollaire (page 54).} \emph{Quand on connaît les signes pour une
face, on les connaît pour toutes.} La page l'appuie sur « la connexité de
$\mathcal{X}$ », qui n'est démontrée nulle part dans le dossier : elle revient à
dire qu'on passe de toute position non supercritique à toute autre par une suite
de spécialisations et de générisations.\footnote{C'est la connexité du graphe des
positions relatives ; dans le langage des matroïdes orientés, celle de l'espace
des extensions en rang $3$ (voir la note de la première section). La
démonstration du dossier suppose en outre, il le dit, que ni $F$ ni $F'$ n'est
une face exceptionnelle.}

La démonstration suppose que ni $F$ ni $F'$ ne correspond à une position
supercritique, et « il faudrait que j'examine séparément le cas de telles
faces ». Si cela marche, on aurait une description simple des signes à partir de
la seule structure $(\mathcal{X}, K, S_K)$, sans points critiques et sans le
graphe $L$.

\pagerange{55}{63}
\section*{La monodromie diédrale (pages 55 à 63, 3 et 4 janvier 1984)}

\subsection*{Le transport le long d'une chaîne de faces}

Une page de calculs au crayon (page 55) met la loi de symétrie en formules.
Soit une chaîne de faces $F_0, F_1, \ldots, F_p$, chacune recollée à la suivante
le long d'une arête $a_i$. Identifions chaque $\operatorname{Pol}(F_i)$ à
$\mathbb{Z}/2n\mathbb{Z}$ par une bijection $\varphi_i$ qui respecte la
structure polygonale. Le transport $\sigma_{a_i} : \operatorname{Pol}(F_i) \to
\operatorname{Pol}(F_{i+1})$ se lit alors comme une application
$\mathbb{Z}/2n\mathbb{Z} \to \mathbb{Z}/2n\mathbb{Z}$ qui respecte la structure
polygonale, renverse l'orientation et envoie un point donné $l_i$ sur $0$ :
c'est la réflexion
\[
\sigma_i(t) = l_i - t .
\]
Le composé de $k$ réflexions consécutives est
\[
\rho_k(t) = (\sigma_{k-1} \cdots \sigma_1 \sigma_0)(t) = (-1)^{k}\bigl(t -
\chi_k\bigr), \qquad \chi_k = l_0 - l_1 + l_2 - \cdots + (-1)^{k-1} l_{k-1} .
\]
C'est la formule de la page 55.\footnote{La page indexe la chaîne par la lettre $n$, qui désigne ailleurs le
nombre des pseudo-droites ; nous l'appelons $k$. Une première valeur,
« $(-1)^{p-1}(\chi_p - t)$ », est biffée ; la formule gardée est la bonne,
comme on le vérifie sur $\sigma_1\sigma_0(t) = l_1 - l_0 + t$ et
$\sigma_2\sigma_1\sigma_0(t) = l_2 - l_1 + l_0 - t$, que la page écrit aussi.}

Pour le tour d'une position supercritique $D_i$, la chaîne a $p = 2n$ faces, le
signe est $+1$, et la somme alternée $\chi_{2n}$ vaut
\[
\chi_{2n} = \sum_{s \in D_i \cap S} \bigl( \nu(s) - 1 \bigr) = n - 1 ,
\]
puisque chaque autre pseudo-droite coupe $D_i$ une fois, en un sommet de
multiplicité $\nu(s)$ où passent $\nu(s) - 1$ d'entre elles.\footnote{Les
nombres $l_s = n - \nu_s$ et $l_a = (n - 1) - (\nu_s - 1) - (\nu_t - 1)$, dont la
différence est $\nu_t - 1$, sont ceux que la page écrit pour les sommets et les
arêtes de la chaîne ; c'est ce qui fait de la somme alternée une somme des
$\nu - 1$. L'identification précise de la chaîne ($2n$ faces, alternativement
de type sommet et de type arête) est la sienne, et nous ne l'avons pas refaite.}

\subsection*{Les images de $\pi_1$ (3 janvier)}

D'où la formule encadrée de la page 57 : \emph{le transport parallèle le long
d'un demi-tour autour d'une position supercritique est la translation}
\[
\lambda(t) = t - (n - 1)
\]
\emph{de $\mathbb{Z}/2n\mathbb{Z}$.} Or $\operatorname{pgcd}(n - 1, 2n) =
\operatorname{pgcd}(n - 1, 2)$ vaut $1$ si $n$ est pair et $2$ si $n$ est
impair. Donc l'image de $\pi_1(\mathcal{X}) \to \mathbb{D}_{2n}$ contient, si
$n$ est pair, tout le groupe des rotations $\mathbb{D}^{+}_{2n} \simeq
\mathbb{Z}/2n\mathbb{Z}$, et elle est d'indice $1$ ou $2$ ; si $n$ est impair,
elle contient le sous-groupe des rotations d'un nombre pair de pas, isomorphe à
$\mathbb{Z}/n\mathbb{Z}$, et elle est d'indice $1$, $2$ ou $4$.\footnote{La page
écrit, pour $n$ pair, « contient $\mathbb{D}^{+}_{2n} \simeq
\mathbb{Z}/n\mathbb{Z}$ » ; le groupe des rotations d'un polygone d'ordre $2n$
est $\mathbb{Z}/2n\mathbb{Z}$, et c'est bien lui qu'engendre une translation
d'un nombre de pas premier à $2n$. Pour $n$ impair, sa conclusion (« indice
$1$, $2$ ou $4$ ») est la bonne.}

Sur le déployé $\widetilde{\mathcal{X}}$ le long des composantes supercritiques
de $K$, où l'on ne franchit plus les positions supercritiques, seuls les tours
complets restent, et l'on trouve la translation $t \mapsto t - 2(n - 1)$ ;
comme $\operatorname{pgcd}(2(n-1), 2n) = 2$, l'image contient
$\mathbb{Z}/n\mathbb{Z}$, sous-groupe d'indice $4$ de quotient
$\mathbb{Z}/2 \times \mathbb{Z}/2$. Cela ne dit rien de l'image dans le groupe
$\lbrace \pm 1 \rbrace$ par le caractère d'orientation $\mathbb{D}_{2n} \to
\lbrace \pm 1 \rbrace$.

\subsection*{La parité des pas}

Pour atteindre ce caractère, il faut un lacet qui retourne l'orientation, et
c'est une affaire de parité : chaque pas élémentaire est une réflexion. En
faisant un demi-tour qui ramène $(D, \omega)$ à $(D, -\omega)$, on compte
\[
2\,(\text{sommets balayés}) + 2\,(\text{droites par-dessus lesquelles on
balaie}) + 1 \ (\text{décollage}) + 1 \ (\text{atterrissage}),
\]
« merveilleux, c'est toujours pair ! ». Mais le raisonnement est incomplet : il
faut compter dans le revêtement où l'on traverse les bandes supercritiques, et
elles sont d'ordre $2(n-1)$, non $2n$. En partant plutôt d'une position voisine
de $D_i$, passant par un sommet $s$ et déduite de $D_i$ par un petit pivotement
autour de $s$, chaque droite franchie compte une seule fois, et le nombre de pas
est $2\,(\ldots) + \nu(s)$. L'image du circuit dans $\lbrace \pm 1 \rbrace$
dépend donc de la parité de l'ordre $\nu(s)$ de $s$, et dès que $\Sigma$ a un
sommet double, $\pi_1(\widetilde{\mathcal{X}}) \to \lbrace \pm 1 \rbrace$ est
surjectif. Pour $n$ pair, joint à ce qui précède, $\pi_1(\mathcal{X}) \to
\mathbb{D}_{2n}$ est surjectif ; « mais peut-être pas »
$\pi_1(\widetilde{\mathcal{X}}) \to \mathbb{D}_{2n}$ — la question est tranchée
à la page 86.\footnote{Pages 57-58, les deux revêtements comparés ne se
distinguant sur la page que par un petit signe peu net au-dessus du tilde. La
parenthèse du nombre de pas contient deux mots illisibles.}

Un tour complet, carré de deux demi-tours, donne un transport qui conserve
l'orientation ; la règle, au bout du compte, est que balayer au-dessus d'un
sommet compte pour deux pas, et au-dessus d'une position supercritique pour un.
« Voilà enfin une façon de voir que $\pi_1(\mathcal{X}) \to \mathbb{D}_{2n}$ est
surjectif » : partant d'une position stable $D$ proche d'une position
supercritique $D_0$, on la spécialise vers le sommet $s$ le plus proche, on fait
un demi-tour autour de $s$ et l'on revient ; c'est un nombre pair d'opérations
qui change l'orientation de $D$ et fixe $s$ — c'est la réflexion par rapport à
$s$. Pour $n$ pair, cette réflexion et les rotations engendrent
$\mathbb{D}_{2n}$. Pour $n$ impair, elles n'engendrent que le sous-groupe
d'indice $2$ formé des rotations d'un nombre pair de pas et des réflexions par
rapport aux sommets ; il faut encore une réflexion par rapport au milieu d'un
côté, que fournissent les demi-circuits de la page 63 et la page
85.\footnote{Pages 59-60. La surjectivité « enfin » obtenue ne vaut donc telle
quelle que pour $n$ pair ; c'est le complément que nous indiquons. Les pages 57
à 60 sont barrées de longs traits obliques qui ne masquent pas le texte, sans
que rien n'indique qu'ils l'annulent.}

\subsection*{Les caractères (4 janvier)}

Pour tout chemin combinatoire $c$ de faces, de $F = F_D$ à $F' = F_{D'}$, on a
un isomorphisme de transport parallèle $\operatorname{Pol}(c) :
\operatorname{Pol}(F) \to \operatorname{Pol}(F')$, composé des réflexions le long
des arêtes traversées, et donc pour un circuit une représentation
\[
u_D : \pi_1(\mathcal{X}, F_D) \longrightarrow \operatorname{Aut}\bigl(
\operatorname{Pol}(D) \bigr) \simeq \mathbb{D}_{2n} .
\]
Le groupe $\mathbb{D}_{2n}$, pour un polygone d'ordre pair, a pour abélianisé
$\mathbb{Z}/2 \times \mathbb{Z}/2$, et donc trois caractères non triviaux à
valeurs dans $\lbrace \pm 1 \rbrace$ : l'effet sur les orientations
$\chi_{\ast}$ ; l'effet sur les deux $n$-gones inscrits dans le polygone (les
sommets de rang pair, ceux de rang impair), $\chi_s$ ; et l'effet sur les deux
classes d'arêtes, $\chi_a$. Comme $\omega(F) \simeq \varepsilon_s(F) \wedge
\varepsilon_a(F)$, où $\varepsilon_s$ et $\varepsilon_a$ sont les ensembles à
deux éléments des $n$-gones inscrits et des classes d'arêtes,
\[
\chi_{\ast} = \chi_s\,\chi_a, \qquad \text{ou, plus symétriquement,} \qquad
\chi_{\ast}\,\chi_s\,\chi_a = 1 .
\]
Ce sont les relations de la page 62.\footnote{La page écrit le second facteur
de la seconde formule $\rho_s$ ; c'est manifestement $\chi_s$. Les « $n$-gones
circonscrits » de la page 60, lecture douteuse, sont les deux classes
d'arêtes.} Deux autres caractères s'y ajoutent : le caractère d'orientation $\chi_{\mathcal{X}}$ de la surface, et
le caractère « combinatoire » $\chi_c$, qui donne la parité de la longueur d'un
lacet. Comme le système local des orientations de la surface est défini par des
transitions \emph{opposées} au transport parallèle, $\omega_{\mathcal{X}}(c) =
\operatorname{par}(c)\,\omega_{\ast}(c)$, d'où
\[
\chi_{\mathcal{X}} = \chi_{\ast}\,\chi_c .
\]
C'est ce que dit la page 62, à une lettre près.\footnote{La formule encadrée de
la page 62 porte $\chi_{\mathcal{X}} = \chi_a\,\chi_c$, l'indice étant un petit
rond ; la formule qui la précède donne $\chi_{\ast}\chi_c$, comme la dernière
ligne de la même page (« et surtout $\chi_{\mathcal{X}} = \chi_{\ast}\chi_c$ »),
et c'est ce que nous écrivons.} On choisit pour caractères fondamentaux $\chi_{\ast}$, $\chi_s$, $\chi_c$, et
l'on a $\chi_a = \chi_{\ast}\chi_s$, $\chi_{\mathcal{X}} = \chi_{\ast}\chi_c$.

Sur les circuits fondamentaux, les valeurs sont les suivantes. Un circuit
$l_{D_i}$ qui franchit une fois la composante exceptionnelle de $K$ attachée à
$D_i$ agit par $u(l_{D_i}) = $ translation de $n - 1$, avec
$\chi_{\mathcal{X}}(l_{D_i}) = 1$ et $\chi_s(l_{D_i}) = (-1)^{n-1}$. Un
demi-circuit autour d'un point de $X$ — dans une face, sur une arête ou en un
sommet — agit par une symétrie ; on les obtient toutes, et donc
\[
u_D : \pi_1(\mathcal{X}, F_D) \longrightarrow \operatorname{Aut}(F_D)
\ \text{ est surjectif — « ouf ! ».}
\]
Enfin $\chi_c$ est \emph{trivial} : si $\operatorname{sing}_D$ désigne le
nombre de sommets de $\Sigma$ sur $D$, le poids d'un pas entre deux positions
voisines est congru à $\operatorname{sing}_{D'} - \operatorname{sing}_D$ modulo
$2$ — on se ramène à un pas élémentaire, où c'est immédiat —, donc le poids d'un
circuit est pair. D'où
\[
\chi_{\mathcal{X}} = \chi_{\ast} :
\]
« il n'y a pas d'autre caractère que ceux provenant de $\mathbb{D}_2$ ! ».
L'orientabilité de la surface est lue sur l'orientation des positions
relatives.\footnote{Page 63, marquée d'un astérisque et introduite par « Il
semble que ». La page de calculs 81, barrée, donne à un caractère de parité une
valeur non triviale sur les demi-tours autour d'un point d'ordre impair ; voir
la section des pages 81 à 100.}

\pagerange{64}{80}
\section*{Huit recollements, et une erreur de deux semaines (pages 64 à 80, 3 janvier 1984)}

\subsection*{Les huit façons}

Le choix du recollement des pages 37-38 est repris de zéro. Pour recoller
$\Delta_D$ à $\Delta_{D'}$ le long d'une transition, il y a a priori huit façons
« standard » : identifier l'un des deux arcs de transition $a$, $a'$ de
$\Delta_{D'}$ (l'arc et son antipodique) à l'un des deux arcs $b$, $b'$ de
$\Delta_D$, de l'une des deux manières. Toutes donnent le même graphe
$\Gamma$, mais un choix et celui qu'on en déduit en changeant partout
d'antipodique définissent la même semi-circularisation et les mêmes transports
d'orientation, donc des surfaces canoniquement isomorphes : il n'y a que quatre
surfaces différentes, et l'on peut supposer que l'arc de $\Delta_D$ est
$b$.\footnote{Page 64 ; les troisièmes termes des triplets, lus $\rho$ et
$\sigma'$, sont peu nets, et le bas de la page, puis le haut de la page 65, sont
barrés.}

Si l'on veut une application naturelle de la surface recollée dans $X$, qui
sur $\Delta_D$ et $\Delta_{D'}$ se réduise aux plongements, il faut
recoller $b$ à $a$ ; mais en remplaçant $\Delta_D$ par un disque un peu plus
petit $\Delta^{*}_D$, dont le bord passe par le sommet $s$, et en recollant le
long de l'arc $b^{*}$ correspondant, on obtient une application dans $X$ dans
tous les cas. Le choix le plus naturel pour l'application dans $X$ semble alors
$b \leftrightarrow a'$ : la rive d'arrivée « fait face » à celle de départ, une
rive se continue par la suivante, et $\mathcal{X} \to X$ est une
\emph{immersion}. Mais avec ce choix on n'a plus d'interprétation des trous, ou
des sommets quand on les a bouchés, par la géométrie de $\Sigma$.

\subsection*{« Je m'étais mélangé les pinceaux deux fois »}

En vérifiant le transport des orientations, Grothendieck s'aperçoit que, dans
les deux cas, l'orientation de $D'$ qu'on obtient en continuant une orientation
de $\Delta_D$ à travers l'arc recollé est \emph{opposée} à celle que donne le
transport parallèle. Et il fait le compte de ses deux dernières semaines :
\begin{quote}
Pendant ces deux dernières semaines, c'est cette opération de recollage que
j'ai (\ldots). (\ldots) je m'étais mélangé les pinceaux \emph{deux} fois, il me
semble : 1) je m'étais trompé sur les isom.\ de transition pour les
orientations, et 2) je m'étais trompé également sur l'identification associée à
une op.\ de générisation pour $D'$ par un arc sur $\partial\Delta_{D'}$ ; on
prenait au lieu de la rive « d'en face » celle qui se déduit de $a$ par transport
parallèle, alors que c'est l'autre (l'antipodique), qu'il fallait
prendre.\footnote{Page 67 ; les points de suspension entre parenthèses marquent des mots illisibles, et
« d'en face » est une lecture douteuse.}
\end{quote}

\subsection*{Les quatre possibilités}

Le bilan tient en une liste.
\begin{itemize}
\item[A)] \emph{Recollements qui donnent une application naturelle
$\mathcal{X} \to X$} : l'identification des arcs est celle du transport
parallèle, à l'antipodisme près, et, sur les orientations, la transition
$\omega(D) \to \omega(D')$ est l'\emph{opposé} du transport parallèle.
$A_1$ : on recolle l'arc de spécialisation $b$ de $\Delta_D$ à l'arc $a'$ de la
rive opposée ; $\mathcal{X} \to X$ est une immersion. $A_2$ : on recolle $b$ à
l'arc $a$ qui s'en déduit par transport parallèle ; $\mathcal{X} \to X$ est une
immersion sur chaque face, mais fait un « retournement » (un pli) le long de
chaque arc de transition.
\item[B)] \emph{Recollements sans application naturelle dans $X$} : la
transition sur les orientations \emph{est} le transport parallèle. $B_1$ : on
repère la générisation par l'arc qui fait face — un cas qu'il avait écarté
« un peu vite ». $B_2$ : par l'arc déduit par transport parallèle — un cas où il
est impossible de faire tourner une position autour d'un trou.
\end{itemize}
Vérification faite, « sauf erreur », $A_1$, $A_2$ et $B_1$ donnent les bons
circuits autour d'une position supercritique, et seul $B_2$ les dérange
complètement — « or il semblerait que c'était sur $B_2$ que je m'étais arrêté
pendant une paire de semaines, Dieu sait pourquoi ! ». La plus intéressante
après $A_2$ serait $A_1$, mais il n'arrive pas à en comprendre les
circuits.\footnote{Pages 68 à 70. Un alinéa barré de la page 69, poursuivi en
haut de la page 70, attribuait l'ennui aux arcs de spécialisation
exceptionnels ; « bâton-dans-les-roues », page 69, est une lecture hasardeuse.}

\subsection*{La surface des positions stables}

On revient (page 71) à une surface faite des seules positions \emph{stables}.
Pour une telle $D$, $\Delta_D$ n'a que des arcs de spécialisation, et deux
positions $D$, $D'$ se déduisent l'une de l'autre par balayage au-dessus d'un
sommet ; on recolle le long des antipodiques des arcs de spécialisation, et il
n'y a plus d'ambiguïté sur les arcs. Pour une figure symétrique, les couples
$(a, b')$ et $(b, a')$ sont exclus par la symétrie, et $(a', b')$ se ramène à
$(a, b)$ par antipodisme : restent deux choix du recollement le long de $(a,
b)$, « de prolongation des rives ».
\begin{itemize}
\item Recoller par le transport parallèle : on obtient une immersion $\varphi :
\mathcal{X} \to X$, un graphe $K$ dont les faces correspondent aux positions
stables, et $L = \varphi^{-1}(X_1)$ qui matérialise cette correspondance, à
condition de renfler chaque face en ses points critiques. Mais l'orientation
transportée par le recollement est, « avec surprise », l'opposée de celle du
transport parallèle, et les circuits ne sont pas les bons autour des positions
supercritiques.
\item Recoller par l'opposé du transport parallèle : c'est alors la transition
sur les orientations qui est le transport parallèle, et, contrairement à ce
qu'il croyait, on trouve les bons circuits autour des positions supercritiques.
« Mais les circuits plus généraux restent mystérieux pour moi. »
\end{itemize}
Les deux procédés donnent la même structure polygonale sur les arcs de
spécialisation. Il y en a une autre : sur $D$ elle-même, et non plus sur
$\widetilde{D}$, les arcs de spécialisation sont deux à deux disjoints, ce qui
donne un polygone d'ordre $n$, réalisé par un disque $n$-gonal. La surface
faite de $2n$-gones ne serait-elle pas un revêtement double de celle faite de
$n$-gones, ramifié aux centres des faces ? Les antipodismes des $\Delta^{*}_D$
sont compatibles aux identifications, et le quotient devrait être la carte
envisagée.\footnote{Pages 74-75, dont plusieurs mots sont d'une lecture
incertaine ; l'idée d'obtenir « simultanément, des deux façons antipodiques »
les disques, pour faire un ruban de Möbius, accompagne cette question.}

La description combinatoire de la surface des positions stables tient en trois
données :
\begin{enumerate}
\item le graphe $\Gamma$, dont les sommets sont les positions stables et les
arêtes orientées les couples $(D, a)$, $a$ un arc de transition sur $D$ — un
arc de spécialisation ou l'antipodique d'un tel —, l'opposée de $(D, a)$ étant
$(D', a')$, où $D'$ se déduit de $D$ par balayage à travers l'antipodique de
$a$ et $a'$ est l'antipodique du transporté de $a$ ;
\item la structure polygonale évidente sur les arêtes issues de $D$ ;
\item le transport des orientations, inverse du transport parallèle.
\end{enumerate}
« Hélas on ne trouve pas les bons circuits supersinguliers » ; en prenant pour
3) le transport parallèle, on les trouve, mais on perd l'application naturelle
$\mathcal{X} \to X$. Une variante qui recolle $a$ à $b'$ et $b$ à $a'$ garde une
application $\mathcal{X} \to X$, non plus une immersion mais une application
pliée le long des arêtes de $K$ ; elle a la même qualité par antipodisme et les
mêmes suites de positions dans les circuits, donc le même défaut.

\subsection*{Le cas des vraies droites}

Que se passe-t-il quand $\Sigma$ est un système de \emph{droites} d'un vrai plan
projectif réel $X$ ? À chaque droite $D_i$ correspond un point $\delta_i$ du
plan dual $X^{\vee}$, et à chaque sommet $D_i \cap D_j$ la droite
$\delta_i\delta_j$ de $X^{\vee}$. Supposons $\Sigma$ simple — les $\delta_i$ non
alignés trois à trois — et, quitte à perturber les $\delta_i$, les sommets de
l'arrangement des $n(n-1)/2$ droites $\delta_i\delta_j$ doubles en dehors des
$\delta_i$ eux-mêmes. Une droite $D$ change de position relative exactement quand
elle passe par un sommet $D_i \cap D_j$, c'est-à-dire quand son point dual
traverse la droite $\delta_i\delta_j$ : les positions relatives de vraies
droites sont les cellules de cet arrangement dual, les faces étant les
positions stables et les points $\delta_i$ les positions supercritiques.

Soit $F$ une face, $D = D_F$ sa position, et $a_\alpha$ ses arêtes, portées par
les droites $\Delta_\alpha$ duales des sommets $s_\alpha$ de $\Sigma$ ; chacune
correspond à une spécialisation $D \to D_\alpha$ vers $s_\alpha$, sommet d'un
triangle de spécialisation de $D$. Les orientations de $F$ correspondent à
celles de $D$, et le transport continu d'une orientation de $F$ vers la face
voisine le long de $a_\alpha$ correspond au transport continu — « parallèle » —
de l'orientation de $D$ au cours du balayage. Ce dessin « paraît rendre
plausible le choix naïf » du transport d'orientation, celui qui donne les bons
circuits ; et l'absence d'application naturelle $\mathcal{X} \to X$ « se
rapproche du fait qu'il n'y a pas d'application naturelle $X^{\vee} \to X$ ».

Suivre une arête $a_\alpha$ revient à faire pivoter $D$ autour d'un point $x$
jusqu'à ce qu'elle passe par $s_\alpha$, puis à la faire pivoter autour de
$s_\alpha$ ; l'orientation de $D$, celle du point dual sur $\Delta_\alpha$ et la
rive de balayage forment alors un système « anti-compatible » : l'orientation
de la position est opposée à celle qu'induit l'orientation de $X$ en $s_\alpha$
sur la rive balayée. On pivote ainsi jusqu'au premier sommet rencontré, qui est
aussi celui où le point dual atteint le sommet $\Delta_\alpha \cap \Delta_\beta$
de la face, et l'on continue.\footnote{Pages 77 à 79 ; « antiscompatible » est
une lecture probable, et plusieurs mots de la fin de la page 79 se lisent mal.
La page 80 est une page de dessins : des demi-circuits autour de points d'ordre
$1$, $2$ ou $0$, un faisceau de lignes grises allant d'un point du bord d'un
disque au point antipodique. Les autres figures des pages 64 à 76 montrent les
disques $\Delta_D$, $\Delta_{D'}$ et leurs arcs $a$, $a'$, $b$, $b'$, le disque
rétréci $\Delta^{*}_D$ (le « gribouillis » de la page 65), et les orientations
$\omega$, $\omega'$ ; la page 72 est une page de calculs au crayon, $N = 4n$
involutions $\sigma_0, \ldots, \sigma_{N-1}$ dont deux consécutives
commutent.}

\pagerange{81}{86}
\section*{Surjectivité, et une généralisation (pages 81 à 100, 5 janvier 1984)}

\subsection*{Une table barrée}

La page 81, barrée de trois longs traits, dresse la table de trois caractères
de la monodromie : $\chi_0$, l'orientation de la position ; $\chi_1$, l'effet
sur les deux $n$-gones inscrits ; $\chi_2$, la parité. Le caractère
d'orientabilité de la surface y est $\chi_0\chi_2$, ce qui est la relation
$\chi_{\mathcal{X}} = \chi_{\ast}\chi_c$ de la page 62 sous d'autres noms. Mais
la table donne à $\chi_2$, sur un demi-tour autour d'un point d'ordre $\alpha$,
la valeur $\alpha$ (modulo $2$), non triviale pour $\alpha$ impair, alors que la
page 63 affirmait $\chi_c$ trivial. Les deux pages ne portent peut-être pas sur
la même surface — celle de la page 85 permet de franchir les positions
supercritiques — ; la page 81 est barrée, et le dossier ne revient pas sur la
question.\footnote{La page 81 note les valeurs additivement ($0$, $1$) et
multiplicativement à la fois ; les lettres lues $\xi$ et $\eta$ dans ses
tableaux sont d'une lecture incertaine. Elle conclut : pour $n$ pair,
« $\pi \to \mathbb{D}_{2n}$ surjectif ». La page 83 porte un compte, $n - \nu -
(\nu - 1) - (\nu' - 1) + \nu = n - (\nu' - 1)$, et les pages 82 et 84 des
dessins de bandes en zigzag qui balaient un disque de pseudo-droites.}

\subsection*{Le revêtement principal est connexe}

Le 5 janvier, le résultat. Soit $\mathcal{X}$ la surface des positions non
supercritiques, sans bord, où l'on permet de franchir les positions
supercritiques par une arête ou par un sommet, compatiblement aux poids. Alors
$\pi_1(\mathcal{X}, D) \to \operatorname{Aut}\operatorname{Pol}(D) \simeq
\mathbb{D}_{2n}$ est surjectif : son image contient, pour chaque arête et chaque
sommet, la symétrie correspondante, réalisée par un lacet qui fait faire à $D$
un demi-tour autour d'un sommet.

Si l'on interdit les franchissements — c'est la surface
$\widetilde{\mathcal{X}}$ —, le même raisonnement ne donne que les réflexions
par rapport aux \emph{côtés}, pas celles par rapport aux sommets. Elles
engendrent le sous-groupe d'indice $2$ formé des rotations d'un nombre pair de
pas et des $n$ réflexions par rapport aux côtés — l'un des trois sous-groupes
d'indice $2$ de $\mathbb{D}_{2n}$, les deux autres étant le groupe des rotations
et celui qu'engendrent les réflexions qui fixent deux sommets.\footnote{La page
85 décrit ce sous-groupe comme formé « par les éléments du sous-groupe des
rotations $\mathbb{D}^{+}_{2n}$ (d'ordre $n$, indice $2$) et les réflexions en
question ». Le groupe des rotations est d'ordre $2n$ ; ce qui est d'ordre $n$
et d'indice $2$ dans les rotations, ce sont les rotations d'un nombre pair de
pas, et c'est ce que nous écrivons. La liste des trois sous-groupes d'indice
$2$ est juste.} L'image de $\pi_1(\widetilde{\mathcal{X}})$ est donc d'indice $1$
ou $2$. En comparant un lacet qui franchit une position supercritique au lacet
qui la contourne par des demi-tours, et en partant d'une position qui passe par
un sommet d'ordre pair, on trouve un lacet de $\widetilde{\mathcal{X}}$ dont
l'effet n'est pas dans ce sous-groupe. Donc, \emph{dans tous les cas,
$\pi_1(\widetilde{\mathcal{X}}) \to \mathbb{D}_{2n}$ est surjectif : le
revêtement principal de groupe $\mathbb{D}_{2n}$ de $\widetilde{\mathcal{X}}$
est connexe.}\footnote{Pages 85-86. La prose de liaison de ces deux pages est
d'une lecture très incertaine, en particulier la distinction entre $n$ pair et
$n$ impair et l'effet sur l'ensemble des polygones circonscrits ; nous ne
donnons que la structure de l'argument et sa conclusion, qui sont lisibles.
C'est la réponse à la question laissée ouverte à la page 58.}

\pagerange{87}{100}
\subsection*{Les positions d'un cercle sur une surface quelconque}

Le même jour, la page 87 généralise le cadre. Soit $X$ une surface compacte
munie d'un graphe $X_1$ qui la décompose cellulairement — le cas qui intéresse
étant celui où tous les sommets de $X_1$ sont d'ordre pair, par exemple un
système de pseudo-droites. Un cercle $D \subset X$ est \emph{en position
standard} s'il rencontre $X_1$, en des points d'ordre pair, et s'il y coupe
toutes les branches de $X_1$. Deux tels cercles sont \emph{en même position
relative} si l'on passe de l'un à l'autre par une isotopie qui garde, à chaque
instant, la position standard et l'intersection avec les sommets : il faut
que $D$ et $D'$ rencontrent les mêmes facettes de $(X, X_1)$, et qu'il existe une
famille continue $f_t$ d'automorphismes de $(X, X_1)$, $f_0 = \mathrm{id}$,
$f_1(D) = D'$. On définit le balayage au-dessus d'un sommet pair, la
spécialisation et la générisation, les arcs de transition sur le revêtement
double $\widetilde{D}$, deux à deux disjoints, et le transport parallèle.

Pour les positions « à une seule rive » — $\widetilde{D}$ connexe —, les
transitions issues de $D$ forment un polygone, et la construction la plus
intrinsèque semble celle où la générisation « tire » sur $D$ et où le transport
des orientations est l'opposé du transport parallèle. Tous les sommets de la
surface sont alors d'ordre $4$, en bijection avec les couples $(D, a)$ où $a$
est un « arc d'orientation » sur $\widetilde{D}$ joignant deux points de $D \cap
S(X_1)$. La situation est la plus symétrique quand $X$ découpé le long de $D$
est un disque — ce qui revient à dire que $X$ est un plan projectif réel et $D$
une pseudo-droite —, condition stable par les opérations
élémentaires.\footnote{Page 89, dont la prose de liaison est très incertaine. La
page 90 est une liste : « a) rives marquées $\Rightarrow$ a') coins marqués ;
b) arcs critiques marqués sur les arêtes ; c) \emph{poids} des faces ». La page
88 ne porte qu'un croquis au crayon.}

\subsection*{Pourquoi deux sortes de sommets seulement}

Les sommets de la surface ne sont pas tous d'ordre $4$ : il y a une exception
quand tout un arc de $\widetilde{D}$ est un arc d'orientation. Ce qui, pour les
pseudo-droites, permet de contrôler les sommets est la propriété suivante, notée
(PS) : \emph{soit $a$ un demi-cercle de $\widetilde{D}$, $D$ position relative
d'un système de pseudo-droites ; alors l'intérieur de $a$ contient un arc de
transition ordinaire, sauf dans deux cas}, l'un où les extrémités de $a$ sont
dans des arcs de générisation au-dessus de sommets de $\Sigma$ et $D$ se déduit
d'une $D_i$ par une petite rotation autour d'un tel sommet (cas où le point est
un sommet), l'autre où elles sont dans des arcs de spécialisation et $D$ se
déduit de même d'une $D_i$ par rapport à une arête (cas où il n'en est
pas).\footnote{Page 92, cas 1) et 2°), marqués en accolade « cas $v \in X_0$ »
et « cas $v \notin X_0$ » ; le texte du second cas est en grande partie
illisible, et une première version (page 91) est biffée.} C'est ce qui assure
qu'il n'y a que deux sortes de sommets : ceux d'ordre $4$, et ceux de type (I),
$(D_i, \beta, \omega)$ — une pseudo-droite, une facette $\beta$ de celle-ci, une
orientation de $X$ en $\beta$. « Sinon, on n'a pas de contrôle, a priori, sur
le type des sommets » : sur une surface quelconque, il n'est pas clair que le
transport parallèle des $\operatorname{Pol}(D)$ provienne d'un système local.

\subsection*{La bande supercritique, et une carte pentagonale}

Pour recoller autour d'une position supercritique, on peut aussi aplatir la bande
supercritique sur son axe. Le transport parallèle qui franchit une position
supercritique n'est alors pas celui que donne l'identification des côtés, mais
son antipodique : pour un chemin sur $\mathcal{X}$,
\[
\tau_X = \underline{a}^{\nu}\,\tau_{\mathcal{X}} ,
\]
où $\tau_{\mathcal{X}}$ est le transport-réflexion sur la surface,
$\underline{a}$ l'antipodisme et $\nu$ le nombre de franchissements d'une
ligne supercritique.\footnote{Page 93 ; les deux dernières lignes de la page se
lisent mal. Le grand dessin de la page montre la bande supercritique, les
droites qui s'y croisent en sabliers, et les comptes « $n - \nu$ » et « $n - 1 -
(\nu - 1) - (\nu' - 1)$ ».}

Une remarque isolée (page 94) : une carte dont les sommets \emph{et} les faces
sont d'ordre pair donne naissance à une suite circulaire de cinq involutions,
$\sigma_0, \sigma_f, \sigma_1, \sigma_s, \sigma_2$, qu'il dit commutantes, et
donc à une carte « pentagonale » dont les sommets sont d'ordre $2$ ou $4$. « Je
n'arrive pas bien à visualiser cette dernière carte, en termes de la carte
initiale\ldots »\footnote{La page ne dit pas lesquelles de ces involutions
commutent ; dans une carte ce sont $\sigma_0$ et $\sigma_2$, qui ne sont pas
voisines dans la suite. Le mot qui qualifie la carte pentagonale est d'une
lecture incertaine.}

La page 96 fixe enfin deux hypothèses : qu'il y ait une pseudo-droite qui ne
passe par aucun des deux sommets $s$, $t$ considérés (on n'est pas dans le cas
$\mathrm{III}_{p,q}$), et une qui ne passe par aucun des points $u$, $u'$ (pas
dans le cas $\mathrm{I}_n$) ; quand il n'y en a qu'une, l'arc de spécialisation
exceptionnel correspondant est de longueur nulle.\footnote{Les pages 95, 97, 99
et 100 sont des pages de dessins en couleur sans texte : des droites coupées
de paires de droites, avec des flèches numérotées ; des disques et des bandes
balayés ; des arbres ramifiés portant de petits disques ; des hémisphères en
perspective. Les dessins de la page 96 marquent des segments de bord de poids
« $\frac12$ ».}

\pagerange{101}{120}
\section*{Rotation autour d'une face, et caractéristique d'Euler (pages 101 à 120)}

Ces vingt feuillets sont, pour plus de la moitié, des dessins ; le texte tient
en quelques pages.\footnote{Pages de dessins sans texte, ou presque : 101 (une
carte en disque et son graphe dual en bleu, des balayages), 102 (des disques de
pseudo-droites balayés), 105, 107, 108 (croquis pâles), 109 (des droites qui
montent en escalier par-dessus des croisements, flèches numérotées), 114 et 115
(des huit et des demi-cercles marqués de signes $\pm 1$, $\varepsilon_1$,
$\varepsilon_2$), 116 (un carré aux côtés I à IV rempli d'arcs), 120 (un ruban
fermé autour d'une droite, un parallélépipède).}

\subsection*{L'opération de rotation}

La page 103 définit une opération sur les triplets $(D, \omega, a)$ — une
position, une orientation, un arc de transition orienté sur $\widetilde{D}$ :
$\rho(D, \omega, a) = (D', \omega', a')$, où
\begin{itemize}
\item[a)] $D'$ se déduit de $D$ en transitant à travers $a$ ;
\item[b)] $\omega'$ est l'orientation \emph{opposée} à celle qu'on déduit de
$\omega$ par transport ;
\item[c)] $a'$ est l'arc de transition de $\widetilde{D}'$ qui \emph{suit}, dans
la direction $\omega'$, le transporté de $a$.
\end{itemize}
C'est la rotation autour d'une face : les triplets $(D, \omega, a)$ sont les
repères de la page 31, et $\rho$ est le parcours $\sigma_1\sigma_0$ du bord
d'une face, pour le transport des orientations opposé au transport
parallèle.\footnote{Le rapprochement avec la page 31 est le nôtre ; il est
exact si l'on prend $\varphi_{\vec{a}}$ égal à l'opposé du transport parallèle,
le choix auquel les pages 64 à 70 ont abouti. La lecture de « rotation » et de
« face », en tête de la page 103, est incertaine.} Il y a huit cas, groupés
selon que $a$ est un arc de spécialisation ou de générisation ; pour trouver
$a'$, on peut soit prendre l'arc de transition qui suit $a$ sur $\widetilde{D}$
dans le sens opposé à $\omega$ et le transporter par balayage ou
anti-balayage, soit transporter d'abord $a$ en $b$ et prendre le successeur de
$b$ pour $\omega'$ — parce qu'une spécialisation induit, entre les arcs de
transition de $\operatorname{Pol}(D)$ non adjacents à $a$ et ceux de
$\operatorname{Pol}(D')$ distincts de l'antipodique de $b$, une bijection. La
page 104 dessine un cycle de cinq positions qui revient à la première ; elle
nomme \emph{arcs d'isotopie} les arcs qui ne sont pas de transition, et note
qu'il faudrait regarder à part les arcs exceptionnels, où il ne devrait pas y
avoir de différence.

\subsection*{La bande supercritique}

Les pages 110 à 113 comptent, sur la bande supercritique, les sommets de part et
d'autre : $n - \nu$ d'un côté, $(n - 1) - (\nu - 1) - (\nu' - 1) = n - (\nu +
\nu' - 1)$ de l'autre, et introduisent un signe $\varepsilon = +1$ si la face
spéciale est « du type sommet », $-1$ sinon — « Comment trouver la valeur de
$\varepsilon$ ? », sans réponse.\footnote{Page 110, feuille en hauteur, couverte
de comptes épars et de dessins de la bande. « face » et « spéciale » sont des
lectures incertaines.} Aux pages 111 et 112, un essai de recollement au voisinage
d'un sommet $s$ fait disparaître le sommet (« \emph{non} ! ») ; un autre le
conserve, mais « je n'arrive pas à suivre les circuits autour d'un trou ». La
page 113 note que, pour deux faces voisines, le nombre $N'$ de sommets d'ordre
$4$ de l'une vaut $N - 1$ ou $N$, si $N$ est celui de l'autre, avec des
inégalités $\alpha', \alpha'' \geqslant 1$, $\gamma \geqslant 2$.

\subsection*{Un brouillon barré}

La page 117, barrée d'un grand X vert, reprend les facettes exceptionnelles du
déploiement : autour de chaque cercle exceptionnel $K(i)$ on dispose d'un
couloir canonique, isomorphe au voisinage tubulaire de $\widetilde{D}_i$ dans
$\widetilde{X}$ ; les faces exceptionnelles adjacentes à $D_i$ restituent le
système $\Sigma$ vu depuis $D_i$ ; et, « à cause de la rectification b) », on ne
peut plus caractériser les droites exceptionnelles de $K$ comme celles qui ne
contiennent pas de sommet de $L'$ — ce sont plutôt celles qui \emph{peuvent} en
contenir plusieurs —, et l'on définit $K'_0$ comme la réunion des composantes du
déployé $K'$ qui n'ont pas de bord global.\footnote{Cette « rectification b) » ne
figure sur aucune page du dossier ; l'alinéa de la page 117 porte la lettre d)
ou c), et la série qui le précède manque. Les deux premières lignes continuent
une phrase commencée ailleurs.}

\subsection*{La caractéristique d'Euler}

La page 119 compte. Si tous les sommets de la surface sont d'ordre $4$, la
somme des ordres des sommets est $4s_0 = 2s_1$, donc $s_1 = 2s_0$, et
\[
\chi(\mathcal{X}) = s_0 - s_1 + s_2 = s_2 - s_0 ,
\]
où $s_2$ est le nombre des positions relatives non supercritiques et $s_0$ celui
des couples $(D, a)$, $a$ un arc d'isotopie du type générisation-générisation,
auxquels s'ajoutent les demi-drapeaux de type $(0, 1)$ de $\Sigma$, au nombre de
$2n(n-1)$ pour un système simple.\footnote{La page écrit $\chi(\mathcal{X}) = 1
- g$ et conclut « $g = 1 + s_0 - s_2$ » : le $g$ qu'elle définit ainsi vaut $1 -
\chi$, et ce n'est pas le genre usuel d'une surface fermée ($\chi = 2 - 2g$ si
elle est orientable, $2 - k$ si elle est la somme de $k$ plans projectifs). Le
« b) » des demi-drapeaux n'est précédé d'aucun « a) » ; le compte $2n(n-1)$ se
vérifie : chacune des $n$ droites porte $n - 1$ sommets, et deux demi-arêtes en
chacun.} C'est le seul endroit où le dossier cherche la topologie de la surface
qu'il construit, et il ne va pas plus loin.

La page 118 esquisse la comparaison des deux transports d'orientation, par
rotation et par le polygone, avec les écritures $\omega_D \wedge \omega_{D'}$ et
$\rho_D \simeq \omega_{D'}$ ; elle se lit trop mal pour être
restituée.\footnote{En tête : « $(\mathcal{X}, K, (\delta_s), (\nu_s), \ldots)$ »,
une liste des données du déploiement ; puis un carré de flèches $\omega(D) \to
\omega(D')$, l'une étiquetée « rot $\alpha$ », l'autre « comb.\ $\beta$ », avec
« $\alpha = -1$ » et « $N$ pair, $\alpha = \beta$ ».}

\pagerange{121}{140}
\section*{Fuseaux et repères (pages 121 à 149)}

Dans les dernières soixante pages, le texte se raréfie et les dessins
dominent. Le point de vue est désormais fixé : la surface
$X^{*}$ des positions\footnote{Les pages 132 à 137 l'écrivent $\widetilde{X}$ ;
nous gardons ce nom pour la sphère, revêtement d'orientation de $X$, et
reprenons celui de la page 150.} a pour \emph{faces} les positions stables, pour
\emph{arêtes} les positions sous-stables, et pour \emph{sommets} les positions
critiques ; dans le cas des vraies droites, c'est exactement l'arrangement dual
de la page 77.\footnote{Pages de dessins sans texte, ou presque, de 121 à 131 :
des disques de pseudo-droites balayés par des faisceaux de positions entre deux
points antipodiques du bord (121, 123, 124) ; des graphes superposés, bleu,
vert et orange, et des « fleurs » de pétales autour d'un sommet (125, 126, 127,
128) ; des disques tangents en huit, marqués $+$ et $-$ (129, 130) ; des disques
dont le bord est numéroté $1, 2, \ldots, 10$, « $11 = 1$ », et une circulation
$\xi_0, \xi_1, \ldots, \xi_{n-1}$ (131). La page 122 porte seulement le mot
« \emph{Non} », souligné, à côté d'un disque où une courbe se recoupe.}

\subsection*{Les fuseaux}

Un \emph{fuseau} est le bigone que balaie une position qui pivote autour d'un
sommet fixé : « Fuseau $\longleftrightarrow$ p.r.\ critique, avec sommet fixé
dessus » (page 132). Il porte trois ensembles à deux éléments, ses deux rives
$a$ (les côtés du bigone), ses deux extrémités $\varpi$ et ses deux orientations
$\Omega$, liés par
\[
\Omega \simeq a \wedge \varpi .
\]
Les positions stables qui entourent un fuseau forment un circuit $D_1, D_2,
\ldots, D_8, D_9 = D_1$ ; ce circuit dépend de la position sous-stable de départ
$D$ avec son sommet $s$, du choix de sa générisation de départ — une rive
$\rho$ de $D$ en $s$ — et d'une orientation $\omega$ de $D$, qui donne l'ordre
dans lequel on parcourt les triangles ; moyennant $\rho$, cela revient à se
donner une orientation $\varpi$ de $X$ en $s$. On parcourt le circuit inverse en
changeant \emph{à la fois} la rive et l'orientation. Donner $(D, \varpi,
\omega)$, c'est donc donner le fuseau avec un « épinglage ».\footnote{Pages 132
et 134 ; les quatre dernières lignes de la page 132 sont d'une lecture très
incertaine, comme le mot « épinglage ». La lettre lue $\varpi$ est un $\omega$
barré ; dans « $\Omega \wedge \varpi$ » elle pourrait être un $\sigma$.}

\subsection*{Les quatre involutions}

Sur l'ensemble $\operatorname{Rep}(X^{*})$ des repères de la surface opèrent
quatre involutions, $\sigma_1$, $\sigma$, $\sigma_2$, $\sigma_0$. $\sigma_2$
fixe $\varpi$ et change la coorientation de $D$ ; $\sigma$ remplace un repère
par son symétrique par rapport au sommet origine, et commute à $\sigma_2$.
$\sigma$ et $\sigma_2$ s'interprètent donc comme des opérations sur les fuseaux,
et il reste à interpréter $\sigma_1$ et $\sigma_0$ : les classes modulo
$\sigma_1$ sont les secteurs angulaires de $X^{*}$, qui correspondent
aux « fuseaux locaux » ; pour $\sigma_0$, la page 137 écrit seulement
« $\sigma_0 = {}$? ».\footnote{L'exposant de $X^{*}$, lu « rep », est
incertain. Au bas de la page 134, isolé : « faces $\leftrightarrow$ fuseaux,
sommets $\leftrightarrow$ droites, arêtes $\leftrightarrow$ bisect.\ ang.\
locaux », qui semble décrire une autre carte, celle dont les faces seraient les
fuseaux. La page 134 dit aussi que deux des quatre opérations, « covariantes »,
commutent.}

Le dictionnaire de la page 135 fixe les cellules.
\begin{itemize}
\item \emph{Faces} de $X^{*}$ : les positions stables ; faces orientées :
les positions stables orientées.
\item \emph{Arêtes} : les positions sous-stables.
\item \emph{Repères} : les quadruplets $(D, s, \omega, \varpi)$, où $D$ est une
position critique, $s \in S \cap D$, $\omega$ une orientation de $D$ et
$\varpi$ une orientation de $X$ en $s$.
\item \emph{Sommets} : les positions critiques ; sommets orientés : les
positions critiques orientées.
\end{itemize}
Un repère est bien un drapeau — un sommet, une arête qui en part, une face qui
contient l'arête — : le sommet est la position critique $D$, l'arête la
position sous-stable qu'on obtient en la faisant pivoter autour de $s$ en
gardant ce seul sommet, et les orientations choisissent la face. La page 137
redit que les repères sont les positions orientées de spécialité au moins $2$,
munies d'un sommet et d'une orientation en ce sommet, c'est-à-dire d'une rive de
$D \smallsetminus \lbrace s \rbrace$.

\subsection*{Le graphe des couples}

La page 138 porte, seul, un graphe en couleur sur les quatre involutions, leurs
produits deux à deux et les couples qu'elles forment : chaque involution et
chaque produit pointe vers le couple auquel il appartient.\footnote{Recomposé en
diagramme d'après le fac-similé par Michel Hua (27 septembre 2026) ; nous le
reproduisons nœud pour nœud d'après la transcription. Sur la page, en orange,
les flèches vers $(\sigma_1, \sigma)$, $(\sigma, \sigma_2)$, $(\sigma_2,
\sigma_0)$ et celles qui partent de $(\sigma, \sigma_2)$ ; en vert clair, les
flèches vers $(\sigma_1, \sigma_2)$, $(\sigma, \sigma_0)$ et, par deux grands
arcs, $(\sigma_1, \sigma_0)$. La flèche issue de $\sigma\sigma_2$ est en tirets.
La page ne dit pas ce que signifient les trois flèches qui joignent des couples
entre eux, et nous ne l'interprétons pas.}

\begin{tikzcd}[column sep=small, row sep=small, nodes={font=\scriptsize}]
 & & & \sigma_1\sigma_0 \arrow[d] & & & \\
 & & & (\sigma_1, \sigma_0) & & & \\
 & & \sigma_1\sigma_2 \arrow[dr] & & & & \\
 & \sigma_1 \arrow[rr] \arrow[d] \arrow[uurr, bend left] & & (\sigma_1, \sigma_2) & & \sigma_2 \arrow[ll] \arrow[d] \arrow[dll] & \\
\sigma_1\sigma \arrow[r] & (\sigma_1, \sigma) \arrow[urr] & & (\sigma, \sigma_2) \arrow[u] \arrow[dr] & & (\sigma_2, \sigma_0) \arrow[dl] & \sigma_0\sigma_2 \arrow[l] \\
 & & \sigma \arrow[ul] \arrow[ur] \arrow[rr] & & (\sigma, \sigma_0) & & \sigma_0 \arrow[ul] \arrow[ll] \arrow[uuuulll, bend right] \\
 & & & \sigma\sigma_2 \arrow[uu, dashed] & \sigma\sigma_0 \arrow[u] & &
\end{tikzcd}

\subsection*{Le plus petit modèle}

La page 139 essaie ces opérations sur un modèle minimal : un segment de $-1$ à
$+1$ passant par $0$, dont les deux moitiés $a$, $b$ ont chacune deux rives, soit
quatre repères $a^{+}$, $a^{-}$, $b^{+}$, $b^{-}$. $\sigma_0$ et $\sigma_2$ y
coïncident et échangent $a^{+}$ et $a^{-}$, $b^{+}$ et $b^{-}$ ; $\sigma_1$
échange $a^{+}$ et $b^{+}$, $a^{-}$ et $b^{-}$ ; $\sigma$ échange $a^{+}$ et
$b^{-}$, $b^{+}$ et $a^{-}$. Posant $\sigma' = \sigma_0 = \sigma_2$, qui commute
à $\sigma$, on a $\sigma_1 = \sigma\sigma'$ : le groupe engendré est un groupe de
Klein qui agit simplement transitivement sur les quatre repères, et la page
note, sous le mot « subdivision barycentrique », les deux manières d'y lire une
carte, $(\sigma_1, \sigma, \sigma_2, \sigma_0) = (\sigma\sigma', \sigma,
\sigma', \sigma')$ et $(\sigma'_0, \sigma'_1, \sigma'_2) = (\sigma\sigma',
\sigma', \sigma)$.\footnote{Le compte du modèle, « $1$ sommet, $1$ face, $2$
arêtes, $4$ repères », a un premier chiffre surchargé, et nous ne savons pas le
rendre cohérent avec le segment dessiné ; les dessins de la page montrent une
sphère avec les points $-1$, $0$, $+1$, $\infty$ et un méridien épais. La page
136 suit un circuit $D_1, \ldots, D_{20}$ autour d'un grand fuseau, les sommets
$s_1, \ldots, s_6 = s_1$, $t_2$, $t_3$, avec un « $D_{20} \overset{?}{=} D_6$ » ;
la page 140 est un grand dessin de pseudo-droites $A, \ldots, G$ dont les points
d'intersection sont nommés $a_0, \ldots, g_{11}$.}

\pagerange{141}{149}
\subsection*{Compter, et balayer}

La page 141 compte, semble-t-il, les positions critiques les plus simples,
celles qui passent par exactement deux sommets d'un système simple. Pour de
vraies droites il y en a une par paire de sommets qui ne sont pas sur une même
droite de $\Sigma$ : il y a $N = n(n-1)/2$ sommets, donc $\frac12 N(N-1)$
paires, dont $n \binom{n-1}{2} = \frac12 n(n-1)(n-2)$ sont sur une même droite,
soit
\[
\frac{1}{2}\,\frac{n(n-1)}{2}\left(\frac{n(n-1)}{2} - 1\right) -
\frac{n(n-1)(n-2)}{2} \;\sim\; \frac{n^4}{8}
\]
positions critiques ; pour des pseudo-droites, rien ne dit que deux sommets ne
portent qu'une position, et c'est la question 1) de la page 153.\footnote{L'interprétation du compte est la nôtre ; la page
ne donne que la formule et, au-dessous, la liste « p.r.\ stables, p.r.\
sous-stables, p.r.\ \emph{critiques} (ou minimales pour la spécialisation) ».}

Les pages 142 à 149 balaient un disque de sept pseudo-droites $A, \ldots, G$ par
une droite mobile, et enregistrent le balayage de deux façons. La page 144
l'écrit comme une suite de permutations,
\begin{gather*}
ABCDEFG,\ ABCEDGF,\ ABCEGDF,\ ABCGEDF, \\
ABGCEDF,\ AGBCEDF,\ \ldots,\ GAEDCFB,
\end{gather*}
chacune se déduisant de la précédente en renversant un ou plusieurs blocs de
lettres voisines ; la page 145 l'écrit comme un diagramme de fils, où chaque
croisement des brins de rangs $i$ et $i+1$ est noté par une transposition
$s_i$ : $s_3, s_4, s_5, s_2, s_6, s_1, s_5, s_3, \ldots$ Ce sont, sous leurs noms
d'aujourd'hui, les \emph{suites admissibles} de Goodman et Pollack et les
\emph{diagrammes de fils} (« wiring diagrams ») de Goodman, où un arrangement de
pseudo-droites se lit comme un mot réduit dans le groupe
symétrique.\footnote{La page écrit les transpositions $\sigma_i$ ; nous les
notons $s_i$ pour ne pas les confondre avec les involutions des repères. La
colonne des générateurs est surchargée et noyée dans une tache d'encre à
partir de la neuvième ligne. Les diagrammes de fils sont dus à Goodman (1980),
les suites admissibles à Goodman et Pollack (1980-1984) ; rien sur ces pages ne
dit que Grothendieck les connaissait. Les pages 142, 143, 146 à 149 sont des
dessins : des arbres de flèches, des cylindres, des disques balayés dont les
sommets sont numérotés $1$ à $19$, « $17 = 1$ ».}

\pagerange{150}{160}
\section*{La surface $X^{*}$ des positions (pages 150 à 160)}

\subsection*{La structure locale}

Soit $X^{*} = \mathrm{Rel}(\Sigma)$ la surface cellulaire des positions
relatives.\footnote{Le nom « Rel » est une lecture douteuse, sous une
surcharge ; l'exposant de $X$ est une étoile. C'est la surface des pages 132 à 137.} Un sommet $\rho$ correspond à une position critique
$D_\rho$ ; soit $S_\rho = D_\rho \cap S$. Une orientation de $X^{*}$ en $\rho$
correspond à une orientation $\varpi_\rho$ de $D_\rho$. Les arêtes qui partent de
$\rho$ correspondent aux couples $(s, \omega)$, $s \in S_\rho$ et $\omega$ une
orientation : faire pivoter $D_\rho$ autour de $s$ dans un sens ou dans l'autre ;
l'arête $a_{s,\omega}$ est la position sous-stable $D_{\rho,s,\omega}$ ainsi
obtenue. Le sommet $\rho$ est donc d'ordre $2\operatorname{card} S_\rho$.

La difficulté n'est pas là. Elle est dans la description de $X^{*}$ au voisinage
d'une \emph{arête} $a_L$, correspondant à une position sous-stable $L$ de
sommet $s$ : on voit bien les faces qui lui sont incidentes, mais pas ses deux
extrémités. Une orientation de $a_L$ revient à une orientation de $X$ en $s$,
donc à une rive de $L \smallsetminus \lbrace s \rbrace$ ; il s'agirait de décrire
l'extrémité de l'arête orientée comme la position critique qu'on obtient en
« bloquant » $L$ par pivotement autour de $s$ dans le sens indiqué. « L'ennui,
c'est qu'il y a plusieurs choix naturels qui viennent à l'esprit, et qu'il
n'est pas sûr qu'aucun de ces (\ldots) ne marche ! »\footnote{Page 150 ; le mot
entre parenthèses est illisible. Pour de vraies droites la réponse est claire —
on pivote autour de $s$ jusqu'à rencontrer un second sommet —, et c'est la
flexibilité des pseudo-droites qui la rend ambiguë ; la remarque est la nôtre.}

\subsection*{Les types de fuseaux}

La page 152 classe les fuseaux en \emph{couloirs}, dont les deux rives sont
supercritiques, \emph{semi-couloirs}, dont une seule l'est, et fuseaux
\emph{ordinaires}. Elle en donne une définition : un fuseau est une composante
connexe $C$ déterminée par deux pseudo-droites $D$, $D'$ (qui peuvent être dans
$\Sigma$) telles que
\begin{itemize}
\item[a)] $\Sigma \cup \lbrace D, D' \rbrace$ soit un système de pseudo-droites ;
\item[b)] $D \neq D'$, et leur point commun soit un sommet de $\Sigma$ ;
\item[c)] $C$ soit un couloir de $\Sigma \cup \lbrace D, D' \rbrace$ ;
\item[d)] $D$ et $D'$ soient critiques pour $\Sigma$.
\end{itemize}
« Chaque fuseau détermine une p.r.\ sous-stable, et inversement » : à isotopie
près, les fuseaux de $X$ correspondent bijectivement aux positions
sous-stables, c'est-à-dire aux arêtes de $X^{*}$ ; $D$ et $D'$ sont les deux
positions critiques qui en sont les extrémités.\footnote{La dernière phrase est
notre lecture de la définition ; elle rejoint la page 132, où le fuseau est
balayé par une position qui pivote autour d'un sommet fixé. L'alinéa b) porte
un mot illisible avant « $= D \cap D'$ ». La \emph{NB} en tête de la page 152 se
lit très mal.}

\subsection*{Le cube des générisations}

Soit $D$ une position qui passe par $\nu = \operatorname{card}(S \cap D)
\geqslant 2$ sommets. En chacun, une générisation choisit de quel côté
décoller $D$ ; ces choix sont indépendants, et
\emph{les générisations stables de $D$ sont les $2^{\nu}$ sommets d'un cube de
dimension $\nu$, les générisations sous-stables en sont les $\nu 2^{\nu - 1}$
arêtes, et l'incidence est la relation de spécialisation.}\footnote{Page 153.
La justification (l'indépendance des choix aux différents sommets) est la
nôtre ; la page donne l'énoncé et les deux nombres. Elle ajoute « Produit des
diagonales d'un polygone convexe (\ldots) $2\nu$ », que nous ne savons pas
interpréter.} La même page pose une question qu'elle laisse ouverte :
\emph{toute position principale — donc critique — qui passe par $S \cap D$
est-elle une spécialisation de $D$ ?}

\subsection*{Les balayages élémentaires}

Les pages 154 et 155 dessinent les balayages élémentaires d'une position entre
deux positions voisines 1 et 2 : $(\mathrm{I}_1)$, autour d'un seul sommet $a$,
asymétrique, la position 1 n'étant pas critique ; $(\mathrm{I}_2)$, autour de
deux sommets $a$, $b$, symétrique, 1 critique si et seulement si 2 l'est ;
$(\mathrm{I}_3)$, autour de trois sommets, asymétrique mais 1 critique si et
seulement si 2 l'est ; et $(\mathrm{I}_n)$ pour $n = 7$, « asymétrique car $n$
impair ».

\subsection*{Les opérations sur un repère}

La page 156 écrit l'action des involutions sur un repère $(D_0, \beta^{*},
\omega^{*})$ — une position, une coorientation, une orientation :
\[
\sigma : (D_0, \beta^{*}, \omega^{*}) \mapsto (D_0, -\beta^{*}, \omega^{*}),
\quad
\sigma_0 : \mapsto (D_0, \beta^{*}, -\omega^{*}),
\quad
\sigma_2 : \mapsto (D_0, -\beta^{*}, -\omega^{*}),
\]
avec $\rho_s = \sigma_2\sigma_1$, $\rho_f = \sigma_1\sigma_0$ et $\rho_s\rho_f =
\sigma$ — les relations de la page 31, où $\sigma = \sigma_2\sigma_0$. Les
orientations $\alpha$ (le long de la droite), $\beta$ (transverse) et $\omega$
(rotation) sont dites « associées » quand $\omega = \alpha \wedge \beta$ ; celles
d'un repère ne le sont pas, $\omega^{*} = -\alpha^{*} \wedge \beta^{*}$. La page
157 compare deux triangles duaux dans deux plans projectifs réels duaux $X$ et
$\check{X}$, et constate que la dualité fait passer d'un système associé à un
système qui ne l'est pas.\footnote{La lecture d'ensemble du texte de la page 157
est très incertaine ; seule la dernière ligne, « $\alpha, \beta, \omega$ associés,
mais $\alpha', \beta', \omega'$ \emph{non associés} », est nette. Les pages 151
et 158 sont des pages de dessins : une vingtaine de disques dont les régions sont
coloriées en deux classes, des rectangles traversés de courbes tressées.}

\subsection*{Les facettes exceptionnelles (1er janvier 1984)}

La page 159, datée du 1er janvier 1984 et barrée d'une croix verte, complète le
déploiement des pages 43 à 46 sur les facettes exceptionnelles. a) $\varphi^{-1}(X_1)
= L \cup K_0$, la correction dont on a parlé. b) Au voisinage des arêtes
exceptionnelles, une facette exceptionnelle a la forme d'un cylindre dont les
bords sont faits de fuseaux ; la loi de symétrie s'étend aux multiplicités des
sommets de $L$ situés sur les deux arêtes ordinaires incidentes. c) Pour une
arête ordinaire qui passe par un sommet exceptionnel, la rive privilégiée est
celle de la face instable. d) Pour avoir la position relative d'une face
exceptionnelle, il faut modifier la construction de l'alinéa (j) de la page 45 en
remplaçant le bord exceptionnel de la face par un chemin qui traverse son
intérieur entre les deux points critiques principaux.\footnote{Date lue
« 1.1.84 », le dernier chiffre pouvant se lire 3 ou 4 ; 1984 est retenu, le
feuillet tombant entre le 31.12.83 et le 2.1.84 (vérifié sur le fac-similé par
Michel Hua le 27 septembre 2026). Les alinéas b) et d) portent de nombreuses
corrections interlinéaires, dont plusieurs illisibles. La page 160 dessine les
facettes cylindriques $F$, $F'$ et les bandes $a\,a'$, $b\,b'$.}

\pagerange{161}{175}
\section*{Dessins, et les axiomes repris (pages 161 à 175)}

Les quinze dernières pages sont presque toutes des dessins en couleur, dans les
conventions des pages précédentes : les pseudo-droites en rouge ou orange, la
position mobile en vert clair, les bandes supercritiques et les couloirs en
vert foncé, les rives marquées en bleu, les triangles en jaune.\footnote{Pages
161 à 164, 166 à 169, 172 à 174 : des disques de pseudo-droites dont le bord est
numéroté $1, \ldots, 5$ et $1', \ldots, 5'$ ; des disques accolés en huit ou en
chaîne ; des étoiles et des croix de droites vertes hérissées de pseudo-droites ;
des faces hachurées reliées par des croisements ; une couronne de pétales autour
d'un cercle ; une sphère coupée de grands cercles et deux polyèdres. La page 165
porte une légende, « Canal : lieu des coins entre bords (bleus) marqués, et
arêtes exceptionnelles instables (aux sommets stables) », lecture en partie
incertaine ; les pages 171 et 173 écrivent des suites circulaires d'étiquettes,
comme « $3\,2\,1\,0\,4\,5\,6\,7\,8\,9\,0'\,1' \ldots 9' \mid 3\,2\,1$ », l'ordre
des pseudo-droites sur $\widetilde{D}$ avant et après un mouvement.}

Une seule page de texte, la page 170, reprend les axiomes du déploiement, cette
fois pour une surface \emph{à bord} $\mathcal{X}$, munie de deux graphes $K$ et
$L$ :
\begin{itemize}
\item[a)] tout sommet intérieur de $K$ est d'ordre $4$, les sommets du bord
d'ordre $1$ ;
\item[b)] on a un déployé $K'$ de $K$ et une section des feuillets des bords de
l'immersion $K' \to \mathcal{X}$ (les rives marquées) ;
\item[c)] les sommets de $L$ sont d'un ordre qui ne se lit pas ;
\item[d)] $L \cap K$ est fini, et en chaque point commun toute branche de $K$
croise toute branche de $L$ ;
\item[e)] chaque arête de $K$ porte \emph{exactement} un sommet de $L$ ;
\item[f)] en chaque sommet intérieur $s$ de $K$, le nombre $\gamma_s$ de points de
$L$ entre $s$ et le sommet de $L$ sur chacune des quatre arêtes ne dépend pas de
l'arête ;
\item[g)] en chaque sommet intérieur $s$, ou bien $s \notin L$ et l'on pose
$\varepsilon_s = 0$, ou bien $s$ est un sommet de $L$ d'ordre réduit $2$ et l'on
pose $\varepsilon_s = 1$ ;
\item[h)] pour toute face $F$ de $K$,
\[
\sum_{s \in S(K) \cap F} \bigl( 2\gamma_s + \varepsilon_s \bigr) \;+\;
\sum_{a \in A(K),\ a \subset F} \gamma_a \;=\; 2n ,
\]
avec $n$ indépendant de $F$.
\end{itemize}
C'est la seconde version des axiomes de la page 43, et non une copie : au
« au plus un sommet de $L$ » par arête de la page 43 répond ici « exactement
un », et les cercles exceptionnels $K_0$, qui portaient les arêtes sans sommet de
$L$, ont disparu au profit du bord.\footnote{Le « exactement » de l'alinéa e)
est suivi d'un ajout illisible, qui semble commencer par « sauf ». Les lettres
$S(K)$, $A(K)$ (sommets et arêtes de $K$) et la condition « $a \subset F$ » sont
des lectures de la transcription. Rien sur la page ne dit explicitement que le
bord remplace $K_0$ ; c'est notre lecture de la différence entre les deux
listes.}

La dernière page du dossier (175) dessine, en tête, le cycle
\[
1 \xrightarrow{\text{spéc}} 2 \xrightarrow{\text{gén}} 3
\xrightarrow{\text{gén}} 4 \xrightarrow{\text{spéc}} 5 = 1 :
\]
le tour d'une face carrée, le même que celui des pages 38 et 39, parcouru à
partir d'un autre côté. Le dossier s'arrête là, sur la cellule élémentaire de la surface qu'il
n'a pas fini de construire.

\end{document}
